% Manuscript prepared for submission to PRX Quantum.
% Two-column APS REVTeX 4.2 article (prxquantum society option).
% Compiles with a standard TeX distribution (TeX Live) or tectonic.
\documentclass[aps,prxquantum,reprint,amsmath,amssymb,floatfix,nofootinbib]{revtex4-2}

\usepackage{mathtools}
\usepackage{amsthm}
\usepackage{bm}
\usepackage{graphicx}
\usepackage{booktabs}
\usepackage{multirow}
\usepackage{enumitem}
\usepackage{microtype}
\usepackage{xcolor}
\usepackage[hidelinks]{hyperref}

\graphicspath{{figures/}}

% Auto-generated macros: every quoted number is defined in
% tables/macros.tex, synced by the package entrypoint cgs-reproduce from
% code/results/macros.tex (built by scripts/make_tables.py from
% results/summary.json).
\input{tables/macros.tex}

% Absorb the few unavoidable few-point overfull boxes caused by long, unbreakable
% typewriter file paths without affecting the spacing of ordinary paragraphs.
\setlength{\emergencystretch}{3em}

\newtheorem{theorem}{Theorem}
\newtheorem{lemma}{Lemma}
\newtheorem{proposition}{Proposition}
\newtheorem{corollary}{Corollary}
\newtheorem{definition}{Definition}
\newtheorem{remark}{Remark}

\DeclareMathOperator{\Tr}{Tr}
% --- Operators used by the error-operator theory section (added) ---
\DeclareMathOperator*{\argmin}{arg\,min}
\DeclareMathOperator*{\argmax}{arg\,max}
\newcommand{\Prob}{\mathbb{P}}
\newcommand{\Expect}{\mathbb{E}}
\newcommand{\Var}{\mathrm{Var}}
\newcommand{\R}{\mathbb{R}}
\newcommand{\C}{\mathbb{C}}
\newcommand{\calH}{\mathcal{H}}
\newcommand{\calG}{\mathcal{G}}
\newcommand{\calO}{\mathcal{O}}
\newcommand{\ii}{\mathrm{i}}
\newcommand{\norm}[1]{\left\lVert #1 \right\rVert}
\newcommand{\abs}[1]{\left\lvert #1 \right\rvert}
\newcommand{\ket}[1]{\left\lvert #1 \right\rangle}
\newcommand{\bra}[1]{\left\langle #1 \right\rvert}

\begin{document}

\title{Zero-Overhead Cancellation of the Leading Trotter Error\\
via Commutator-Graph Scheduling}

\author{Molena Huynh}
\email{molena.huynh@jmp.com}
\affiliation{North Carolina State University, Raleigh, North Carolina 27695, USA}

\date{\today}

\begin{abstract}
A quantum computer is a stochastic device: the Born rule returns one bitstring
per run, so every quantity a simulation reports---fidelities, observables, and
the very error one is trying to control---is a statistical estimate over samples,
and the structure that governs those estimates is the algebra of the operators
being evolved. We introduce \emph{commutator-graph scheduling} (CGS), which
doubles the convergence order of first-order Trotterization at zero gate, qubit,
or depth overhead by making the leading error an explicit, estimable operator on
that algebra. The bottleneck is well known: term order within a Lie--Trotter step
is free, yet ordering heuristics buy only marginal accuracy and no one had
explained why. CGS recasts the leading error as an exact, matrix-free operator, the
\emph{commutator error form}
$E_\pi=\sum_{a<b}[h_{\pi(a)},h_{\pi(b)}]$---a signed sum of Pauli strings over the
edges of the commutator graph, validated against dense exponentials to $10^{-9}$.
We prove that reordering only re-signs $E_\pi$, leaving its norm exactly invariant
on collision-free families, so no fixed ordering escapes the $O(t^2/r)$ rate. The
true free parameter is the per-step \emph{schedule}: a step-alternating
\emph{antithetic} schedule cancels $E_\pi$ and converges at $O(t^3/r^2)$ with the
identical $L\cdot r$ rotations. Across $\Ninstances$ structured Hamiltonians
versus exact statevector references, CGS lifts the empirical rate from
$\SlopeFirst$ to $\SlopeAnti$, reaches fidelity $0.99$ with $\GateSpeedup\times$
fewer rotations than the best first-order ordering, and attains $0.999$ where
first-order Trotter cannot. A matrix-free Krylov reference extends the exact
benchmark to $\MaxN$ qubits, where the order doubling is size independent
($\approx\!2$ vs.\ $\approx\!4$ from $\MinN$ to $\MaxN$ qubits) and the
matched-cost speedup reaches $\SpeedupMaxN\times$. An interpretable scheduler
over ten graph features selects the gate-optimal schedule per instance with
$\LearnedAcc\%$ leave-one-out accuracy. Our contribution is threefold: a
statistical theory of the leading Trotter error (exact ordering invariance,
antithetic cancellation, strict optimality, and a certified shot estimator), a
zero-overhead scheduling algorithm released as the pip-installable package
\texttt{specops-cgs}, and a matched-cost benchmark in which every figure, table,
and quoted number regenerates from a single seeded command.
\end{abstract}

\maketitle

% ============================================================
\section{Introduction}
\label{sec:introduction}

Statistics is not an afterthought in quantum computing; it is structural to it.
Seen from the keyboard, a quantum computer is a stochastic device interrogated
only by sampling: the Born rule means each run returns one bitstring drawn from a
distribution, so everything a simulation ultimately reports is an
estimate---observables carry variance and confidence intervals, and the residual
error of an approximate circuit is itself something one can only \emph{measure} to
finite precision. This viewpoint sharpens the problem of simulating time
evolution. The error one is trying to suppress is a concrete operator on the
Pauli algebra, and controlling it well means representing it exactly, bounding it
rigorously, and estimating its magnitude from shots with a stated risk---a chain
of statistical reasoning that runs from the operator algebra to the sampled
witness. Commutator-graph scheduling, introduced below, follows exactly this
chain: it turns the leading Trotter error into an explicit, matrix-free operator
on a graph, proves how far one can and cannot shrink it, and gives an unbiased
shot-based estimator of its size so the accuracy claim is empirically certifiable
rather than merely asymptotic.

The simulation of quantum time evolution $e^{-\ii H t}$ is central to a wide array
of fields, such as quantum chemistry, condensed-matter physics, and materials
science, and is among the oldest and most promising applications of quantum
computers~\cite{feynman1982simulating,lloyd1996universal,georgescu2014quantum,aspuru2005simulated}.
Its appeal extends from quantum chemistry, where simulated time evolution underlies
phase estimation of molecular
energies~\cite{aspuru2005simulated,mcardle2020quantum}, to the first digital
quantum simulations on trapped-ion and other near-term
hardware~\cite{lanyon2011universal,smith2016manybody}. The workhorse primitive
for this task is the product formula. Writing a Hamiltonian in its Pauli
decomposition,
\begin{equation}
    H=\sum_{j=1}^{L} c_j P_j ,
    \label{eq:ham}
\end{equation}
with real coefficients $c_j$ and Pauli strings $P_j\in\{I,X,Y,Z\}^{\otimes n}$, the
first-order (Lie--Trotter) approximant with $r$ steps and a term ordering $\pi$ is
\begin{equation}
    U_{\mathrm{Trot}}(\pi)=\Bigg(\prod_{j\in\pi} e^{-\ii\,c_j P_j\, t/r}\Bigg)^{\!r}.
    \label{eq:trotter}
\end{equation}
The accuracy of this approximation is governed by the commutators of the summands:
the leading correction to a single step is
$\tfrac{t^2}{2r}\sum_{a<b}[c_a P_a,\, c_b P_b]$, which vanishes for any pair of
terms that commute~\cite{trotter1959product,suzuki1991general,childs2021theory}.
The product-formula construction itself traces to the operator-splitting schemes of
Strang~\cite{strang1968construction} and to Suzuki's fractal
decompositions~\cite{suzuki1990fractal,suzuki1991general}, and it remains
competitive with the more recent quantum-signal-processing and
linear-combination-of-unitaries
primitives~\cite{low2017optimal,low2019qubitization,berry2015taylor,childs2012hamiltonian}
precisely because its cost scales with commutators rather than with a global norm.
The modern, tight analysis of this error in terms of nested commutators provides
the theoretical anchor of the present work~\cite{childs2021theory,childs2018toward}.
This commutator-scaling viewpoint also explains why product formulas are nearly
optimal for geometrically local lattice
Hamiltonians~\cite{childs2019lattice,childs2010limitations}.

A defining feature of the first-order formula is that its implementation cost is one
rotation $e^{-\ii c_j P_j t/r}$ per term per step, so the gate count is $L\cdot r$
\emph{independently} of the ordering $\pi$, while the error is not. The ordering is
therefore a free knob, and a substantial body of empirical work has asked how much
accuracy it can buy~\cite{tranter2019ordering,poulin2015trotter,hastings2015improving,wecker2015solving}.
A parallel line of work has shown that the worst-case Trotter bound is often loose:
errors from successive steps can interfere
destructively~\cite{tran2020destructive,layden2022first}, are bounded by quantum
localization for local observables~\cite{heyl2019quantum,sieberer2019digital}, and
improve markedly on average over input
states~\cite{chen2024averagecase,yi2022spectral}.
The natural structure for this question is the \emph{commutator graph} (or
anticommutation graph), with one node per term and an edge between every pair of
terms that anticommute---the same graph used in measurement grouping for variational
algorithms, where a clique cover into mutually commuting sets reduces the number of
measurement
settings~\cite{verteletskyi2020measurement,gokhale2020optimization,crawford2021efficient,yen2020measuring,jena2019pauli}.
Despite this rich structure, the empirical verdict has been consistently
underwhelming: ordering heuristics move the error only marginally, and their gains
are widely regarded as ``within noise.'' \emph{Why} a free parameter should be so
impotent has not been explained.

We resolve this, and turn it into an asset, by introducing
\emph{commutator-graph scheduling} (CGS). At the center of CGS is the
\emph{commutator error form}, the exact leading-order error operator
\begin{equation}
    E_\pi=\sum_{a<b}\bigl[h_{\pi(a)},h_{\pi(b)}\bigr],\qquad h_j=c_jP_j,
    \label{eq:errorform}
\end{equation}
which we assemble without ever forming a $2^n\times2^n$ matrix, as a signed sum of
Pauli strings indexed by the edges of the commutator graph
(Sec.~\ref{sec:method}). Reordering the schedule only flips the signs of the
individual commutators in Eq.~\eqref{eq:errorform}; it never changes which
commutators appear. The Hilbert--Schmidt norm of $E_\pi$ is therefore \emph{exactly}
ordering invariant whenever the edge commutators land on distinct Pauli strings, and
only a small ``colliding'' fraction is ordering reducible at all
(Theorem~\ref{thm:impotence}). This is the precise, provable content of the
``ordering is within noise'' folklore: no fixed ordering escapes the $O(t^2/r)$
first-order rate.

The resolution is to enlarge the free degree of freedom. The order applied at each
Trotter step need not be the same; the genuinely free object is a per-step
\emph{schedule} $\sigma=(\pi_1,\dots,\pi_r)$, realized at the identical $L\cdot r$
rotation count. Among all schedules, one cancels Eq.~\eqref{eq:errorform}
\emph{exactly}: the \emph{antithetic} schedule, which alternates a term order with
its reverse on successive steps. Each adjacent step-pair is then a palindromic, and
hence symmetric, product formula whose leading $O(t^2/r)$ commutator error vanishes,
upgrading the convergence from $O(t^2/r)$ to $O(t^3/r^2)$ at \emph{no change in gate
count} (Theorem~\ref{thm:antithetic}). The underlying mechanism is the classical
symmetrization of Suzuki~\cite{suzuki1991general,hatano2005finding}; reading it as a
zero-overhead \emph{scheduling} choice---one that the impotence theorem shows a
single fixed ordering can never realize---is new, and it changes the practical
recommendation for first-order simulation.

Our contributions are fourfold.
\begin{enumerate}[leftmargin=*,itemsep=2pt,topsep=3pt,label=(\roman*)]
\item \emph{An exact, estimable error operator.} We introduce the commutator
error form [Eq.~\eqref{eq:errorform}], a matrix-free representation of the
leading Trotter error as a signed Pauli-string sum over the commutator graph
(Definition~\ref{def:errorform}), validated against the dense pair-commutator
sum to $10^{-9}$.
\item \emph{A statistical theory of scheduling, with full proofs.} A single
ordering cannot reduce the leading-error norm below an instance floor that is
exactly ordering invariant on collision-free families
(Theorem~\ref{thm:impotence})---the provable content of the ``ordering is
within noise'' folklore---while the antithetic schedule cancels the leading
error and converges at second order at the identical gate budget
(Theorem~\ref{thm:antithetic}). The accompanying error-operator theory supplies
a quantitative BCH remainder (Theorem~\ref{thm:bchremainder}), strict optimality
of leading cancellation over every fixed ordering (Theorem~\ref{thm:optimality}),
$\mathsf{NP}$-hardness of residual minimization together with a spectral
relaxation (Theorem~\ref{thm:nphard}), and an unbiased shot-based estimator of
the residual with a finite-sample certified confidence interval
(Theorem~\ref{thm:estimator}). The last item makes the statistical viewpoint
structural rather than decorative: a quantum computer is interrogated only
through Born-rule samples, so the accuracy claim itself is delivered as a
certified statistical inference.
\item \emph{A zero-overhead algorithm, released as a pip-installable package.}
CGS combines the antithetic fold with a commuting-clique compression of the
commutator graph (Proposition~\ref{prop:colour}) and an interpretable learned
scheduler over ten listing-order-invariant graph features
(Proposition~\ref{prop:invariance}). The complete pipeline is the open-source
package \texttt{specops-cgs}; its single console command
\texttt{cgs-reproduce} regenerates every figure, table, and quoted number in
this manuscript from one seeded run.
\item \emph{A matched-cost empirical benchmark.} Against exact statevector
references on $\Ninstances$ structured Hamiltonian instances, CGS lifts the
empirical convergence rate from $\SlopeFirst$ to $\SlopeAnti$, reaches fidelity
$0.99$ with $\GateSpeedup\times$ fewer rotations than the best first-order
ordering, and is the only family of schedules to reach $0.999$ within the swept
budget. A matrix-free Krylov reference extends the exact benchmark from the
dense-exponential ceiling to $\MaxN$ qubits; across $n=\MinN$ to $\MaxN$ the
convergence-order doubling (first-order rate $\approx\!2$, antithetic rate
$\approx\!4$) and the matched-cost speedup are size independent
($\SpeedupMaxN\times$ at $n=\MaxN$; Sec.~\ref{sec:scaling}). The learned
scheduler selects the gate-optimal schedule per instance with $\LearnedAcc\%$
leave-one-out accuracy and recovers near-oracle efficiency
($\LearnedMeanGates$ rotations vs.\ the oracle's $\OracleMeanGates$).
\end{enumerate}

The boundaries of the study are explicit. CGS is evaluated by classical
statevector simulation on systems up to $n=\MaxN$ qubits, where an exact reference
is tractable, on synthetic and structured Hamiltonians, under a rotation-count cost
model; the gains are asymptotic in the step size, and we report the crossover
budget. The antithetic schedule is not a new product-formula primitive---the
symmetrization is classical---but its certified, zero-overhead deployment on the
commutator graph, together with the impotence theorem that motivates it and the
learned scheduler that selects it, is new.

This work continues a program on the spectral truncation of operators for
query-efficient quantum optimization and simulation. Earlier entries develop
graph-conditioned trust regions for budgeted parameter
search~\cite{huynh2026gctr,huynh2026gctrcalpro}, the measurement
and query accounting that fixes the cost model~\cite{huynh2026meascost,huynh2026certqb},
and topology-conditioned transfer of optimizer state across
instances~\cite{huynh2026paramxferB}, culminating in
operator-spectral truncated priors and spectral-truncation graph kernels for QAOA
warm starts~\cite{huynh2026ostqaoa,huynh2026topoqaoa}. The present manuscript carries
the same operator-spectral viewpoint from variational optimization to Hamiltonian
simulation: the commutator error form is a matrix-free, edge-resolved truncation of
the leading product-formula error on the commutator graph, and the matched-cost
analysis reuses the query-budget discipline established in that line of work.

% ============================================================
\section{Commutator-graph scheduling}
\label{sec:method}

We now present the CGS framework and its analytic guarantees. We first record the
Pauli algebra that makes an exact, matrix-free simulator and error form possible
(Sec.~\ref{sec:pauli}), then define the commutator error form
(Sec.~\ref{sec:errorform}), prove that ordering is impotent
(Sec.~\ref{sec:impotence_thy}), prove that the antithetic schedule cancels the
leading error at the same gate budget (Sec.~\ref{sec:antithetic_thy}), establish the
commuting-clique compression (Sec.~\ref{sec:clique_thy}), and describe the learned
scheduler that selects among schedules (Sec.~\ref{sec:router_thy}). Throughout, every
result is stated under explicitly named hypotheses with a complete proof, and every
numerical claim is anchored to a quantity that the accompanying code regenerates.

\begin{figure*}[t]
    \centering
    \includegraphics[width=\textwidth]{fig_schematic.pdf}
    \caption{\textbf{Commutator-graph scheduling (CGS).} A structured Hamiltonian
    $H=\sum_j c_j P_j$ is mapped to its commutator graph $\calG(H)$---one node per
    term, an edge between every anticommuting pair, and a greedy coloring into
    mutually commuting cliques. A \emph{schedule} $\sigma$ fixes, for each Trotter
    step, a term order and a fold (first-order, or the step-alternating antithetic
    fold), optionally laying the commuting cliques out contiguously; the realized
    rotation count is matched across schedules. The product formula $U_\sigma$ is
    executed by an exact, matrix-free statevector simulator and scored by fidelity
    against $e^{-\ii Ht}$ and by the rotations needed to reach a target fidelity.
    The leading error operator $E_\pi=\sum_{a<b}[h_{\pi(a)},h_{\pi(b)}]$ is the
    object everything turns on: reordering only re-signs it
    (Theorem~\ref{thm:impotence}), whereas a folded schedule cancels it
    (Theorem~\ref{thm:antithetic}). Three integrity gates underpin the
    benchmark---the fast Pauli-rotation kernel and dense \texttt{expm} agree to
    $10^{-9}$, the matrix-free error form agrees with the dense pair-commutator sum
    to $10^{-9}$, and the matrix-free Krylov reference agrees with dense
    \texttt{expm} to $10^{-9}$---before any number is emitted.}
    \label{fig:schematic}
\end{figure*}

Figure~\ref{fig:schematic} summarizes the pipeline. Each Hamiltonian instance is
decomposed into Pauli terms; the commutator graph $\calG(H)$ places one node per term
and an edge between every anticommuting pair, and a greedy coloring partitions the
terms into mutually commuting cliques. A schedule fixes, for every Trotter step, a
term order and a fold---first-order, or the step-alternating antithetic fold---while
holding the realized rotation count fixed. The product formula is executed by an
exact, matrix-free statevector simulator and scored both by fidelity against
$e^{-\ii Ht}$ and by the rotations needed to reach a target fidelity. We stress that
no number is reported until three integrity gates pass: the fast Pauli-rotation
kernel and dense \texttt{expm} agree to $10^{-9}$, the matrix-free error form agrees
with the dense pair-commutator sum to $10^{-9}$, and the matrix-free Krylov reference
(which extends the exact benchmark to $\MaxN$ qubits) agrees with dense \texttt{expm}
to $10^{-9}$.

\subsection{Pauli algebra and the matrix-free simulator}
\label{sec:pauli}
A Pauli string $P\in\{I,X,Y,Z\}^{\otimes n}$ acts on $\ket{\psi}\in\C^{2^n}$; we
write $\{A,B\}=AB+BA$ and $[A,B]=AB-BA$, and use $X^2=Y^2=Z^2=I$ together with
$\{X,Y\}=\{Y,Z\}=\{Z,X\}=0$. Two structural facts make an exact simulator possible
without ever forming a $2^n\times2^n$ matrix.

\begin{lemma}[Commutation is an $O(n)$ parity test]
\label{lem:parity}
Let $P=\bigotimes_{q=1}^{n}p_q$ and $Q=\bigotimes_{q=1}^{n}q_q$ be Pauli strings.
Call position $q$ \emph{anticommuting} if $p_q\neq q_q$ and neither is $I$, and let
$k=k(P,Q)$ count the anticommuting positions. Then $PQ=(-1)^{k}QP$, so $P$ and $Q$
commute iff $k$ is even. The test runs in $O(n)$ time and forms no $2^n\times2^n$
matrix.
\end{lemma}

\begin{proof}
At each position either $p_q=q_q$ or one is $I$ (then $p_qq_q=q_qp_q$), or $p_q\neq
q_q$ and neither is $I$ (then $p_qq_q=-q_qp_q$). So $p_qq_q=(-1)^{\epsilon_q}q_qp_q$
with $\epsilon_q=1$ exactly on anticommuting positions, and since the tensor product
factorizes across qubits,
$PQ=\bigl(\prod_q(-1)^{\epsilon_q}\bigr)QP=(-1)^{k}QP$ with $k=\sum_q\epsilon_q$.
Hence $PQ=QP$ iff $k$ is even.
\end{proof}

\begin{lemma}[Pauli rotation as a cosine--sine pair]
\label{lem:rotation}
For a non-identity Pauli string $P$ ($P^2=I$) and $\theta\in\R$,
\begin{equation}
    e^{-\ii\theta P}=\cos\theta\,I-\ii\sin\theta\,P,
    \label{eq:rotation}
\end{equation}
so $e^{-\ii\theta P}\ket{\psi}=\cos\theta\,\ket{\psi}-\ii\sin\theta\,P\ket{\psi}$
is one sparse application of $P$ plus a scaled vector add, computable in $O(n\,2^n)$
time without forming the $2^n\times2^n$ matrix.
\end{lemma}

\begin{proof}
Since $P^{2m}=I$ and $P^{2m+1}=P$, regrouping the convergent series of
$e^{-\ii\theta P}$ by parity gives $\cos\theta\,I-\ii\sin\theta\,P$. A single Pauli
string maps each basis state to one basis state with a unit-modulus phase, so
$P\ket{\psi}$ costs $O(n\,2^n)$.
\end{proof}

We emphasize that these identities are not assumed but verified: before any fidelity
is computed, Eq.~\eqref{eq:rotation} is cross-checked against dense
\texttt{scipy.linalg.expm} on random small Paulis to tolerance $10^{-9}$ (the
\texttt{pauli\_algebra\_verified} flag). The product of two Pauli strings is itself a
phased Pauli, $P\,Q=\omega\,R$ with $\omega\in\{\pm1,\pm\ii\}$ and $R$ a Pauli
string, computed in $O(n)$ (the \texttt{pauli\_product} kernel); this is the algebra
on which the commutator error form rests.

\subsection{The commutator error form}
\label{sec:errorform}
The commutator graph $\calG(H)$ has one node per term and an edge $\{j,k\}$ for every
anticommuting pair (Lemma~\ref{lem:parity}). The exact value of the pairwise
commutators that populate the leading error is fixed by the following lemma.

\begin{lemma}[Pauli commutator value]
\label{lem:commvalue}
Let $h_j=c_jP_j$ and $h_k=c_kP_k$ with real $c_j,c_k$. If $P_j,P_k$ commute then
$[h_j,h_k]=0$; if they anticommute then
$[h_j,h_k]=2c_jc_k\,P_jP_k=2c_jc_k\,\omega_{jk}R_{jk}$ with $\omega_{jk}\in\{\pm\ii\}$
and $R_{jk}$ a Hermitian Pauli string, so
$\norm{[h_j,h_k]}_{\mathrm{HS}}=2\abs{c_jc_k}\,2^{n/2}$ and the operator norm is
$2\abs{c_jc_k}$.
\end{lemma}

\begin{proof}
By Lemma~\ref{lem:parity}, $P_kP_j=(-1)^{k(P_j,P_k)}P_jP_k$, so the commutator
vanishes when the terms commute and equals $2c_jc_k\,P_jP_k$ when they anticommute.
The product $P_jP_k=\omega_{jk}R_{jk}$ is a phased Pauli (Sec.~\ref{sec:pauli});
because the commutator of Hermitian operators is anti-Hermitian and $R_{jk}$ is
Hermitian, $\omega_{jk}\in\{\pm\ii\}$. Every Pauli string is unitary with
$\norm{R_{jk}}=1$ and $\norm{R_{jk}}_{\mathrm{HS}}^2=\Tr(R_{jk}^2)=2^n$, which gives
the stated norms.
\end{proof}

\begin{definition}[Commutator error form]
\label{def:errorform}
For an ordering $\pi$, the leading error operator is
$E_\pi=\sum_{a<b}[h_{\pi(a)},h_{\pi(b)}]$. By Lemma~\ref{lem:commvalue} it equals the
edge sum
\begin{equation}
    E_\pi=\sum_{\{j,k\}\in E(\calG)}\varepsilon_{jk}(\pi)\,2c_jc_k\,(P_jP_k),
    \label{eq:edgesum}
\end{equation}
a signed sum of phased Pauli strings indexed by the edges of $\calG(H)$, with
$\varepsilon_{jk}(\pi)=+1$ if $j$ precedes $k$ in $\pi$ and $-1$ otherwise. It is
assembled in $O(L^2 n)$ time without forming a $2^n\times2^n$ matrix, and is
validated against the dense pair-commutator sum to $10^{-9}$ (the
\texttt{error\_form\_verified} flag).
\end{definition}

Equation~\eqref{eq:edgesum} is the certificate at the heart of CGS: it exposes,
exactly and cheaply, precisely how the ordering enters the leading error---through
the signs $\varepsilon_{jk}$, and nothing else.

\subsection{Ordering is impotent}
\label{sec:impotence_thy}
We first record that the gate budget is a schedule invariant for the schedules we
compare, so that any accuracy difference is a difference \emph{at matched cost}.

\begin{proposition}[Rotation-count invariance]
\label{prop:cost}
A first-order schedule with any ordering $\pi$, and the antithetic schedule that
alternates $\pi$ with its reverse, each apply exactly $L\cdot r$ single-Pauli
rotations over $r$ steps. Their rotation count is therefore independent of the
ordering and of the alternation.
\end{proposition}

\begin{proof}
One first-order step is an ordered product of one factor per term, hence $L$
rotations, and $r$ steps give $L\cdot r$. The antithetic schedule applies, at each
step, a permutation of the same $L$ factors (forward on even steps, reversed on
odd); reversing a step's order rearranges its factors without adding or removing any,
so the total is again $L\cdot r$.
\end{proof}

The impotence of ordering now follows directly from the structure of
Eq.~\eqref{eq:edgesum}.

\begin{theorem}[Ordering impotence]
\label{thm:impotence}
Let $H=\sum_j c_jP_j$ with commutator error form $E_\pi$. Reordering acts on $E_\pi$
only by flipping the signs $\varepsilon_{jk}$; the multiset of pair commutators
$\{[h_j,h_k]\}$ is independent of $\pi$. Grouping the edge sum by Pauli string gives
$\norm{E_\pi}_{\mathrm{HS}}^2=2^n\sum_R\abs{\alpha_R(\pi)}^2$ with
$\alpha_R(\pi)=\ii\sum_{\{j,k\}:R_{jk}=R}\varepsilon_{jk}(\pi)\,2c_jc_k\,\widehat\omega_{jk}$
and $\widehat\omega_{jk}=\omega_{jk}/\ii\in\{\pm1\}$. If every edge maps to a
distinct Pauli string (no collisions), then
\begin{equation}
    \norm{E_\pi}_{\mathrm{HS}}^2=2^n\sum_{\{j,k\}}(2c_jc_k)^2
    \label{eq:hsfloor}
\end{equation}
is \emph{independent of $\pi$}. In general only the colliding edges contribute an
ordering-dependent term, so $\norm{E_\pi}_{\mathrm{HS}}$ is bounded below, for
\emph{every} $\pi$, by the ordering-invariant contribution of the non-colliding
edges.
\end{theorem}

\begin{proof}
For an unordered pair $\{j,k\}$, listing $j$ before $k$ contributes $[h_j,h_k]$ and
the reverse contributes $-[h_j,h_k]$; reordering therefore multiplies each
contribution by $\varepsilon_{jk}(\pi)$ and neither adds nor removes a pair. By
Lemma~\ref{lem:commvalue} each anticommuting edge contributes
$\ii\,2c_jc_k\,\widehat\omega_{jk}R_{jk}$. Distinct Pauli strings are
Hilbert--Schmidt orthogonal with squared norm $2^n$, so
$\norm{E_\pi}_{\mathrm{HS}}^2=2^n\sum_R\abs{\alpha_R(\pi)}^2$ with $\alpha_R(\pi)$ as
stated. If a Pauli string $R$ has a single contributing edge, then
$\abs{\alpha_R}=2\abs{c_jc_k}$ regardless of the sign $\varepsilon_{jk}$, so its
contribution is $\pi$-independent and the no-collision case reduces to
Eq.~\eqref{eq:hsfloor}. Edges that are alone on their Pauli string thus contribute a
fixed amount for every $\pi$; the remaining (colliding) edges are the only source of
ordering dependence, and their contribution is non-negative, which yields the lower
bound.
\end{proof}

We emphasize that Theorem~\ref{thm:impotence} is not a statement about any one
heuristic: on a collision-free family the leading-error norm is \emph{identical for
every ordering}, and in general only a small colliding fraction is reducible at all.
No reordering escapes the $O(t^2/r)$ rate. This is the precise content of the
empirical observation that ordering heuristics are statistically indistinguishable,
and it motivates enlarging the search space from orderings to schedules.

\subsection{A free schedule cancels the leading error}
\label{sec:antithetic_thy}
The order applied at each Trotter step is itself free at the same gate budget, and
CGS exploits this with the \emph{antithetic} schedule, which applies a term order
$\pi$ on even steps and its reverse $\pi^{R}$ on odd steps.

\begin{theorem}[Antithetic cancellation]
\label{thm:antithetic}
Fix an ordering $\pi$ and step size $\Delta=t/r$ with $r$ even, and let
$S_\pi(\Delta)=\prod_{a=1}^{L}e^{-\ii\Delta h_{\pi(a)}}$ and
$S_{\pi^R}(\Delta)=\prod_{a=L}^{1}e^{-\ii\Delta h_{\pi(a)}}$ be the forward and
reversed steps. The antithetic step-pair $M(\Delta)=S_{\pi^R}(\Delta)S_\pi(\Delta)$
is symmetric, $M(-\Delta)=M(\Delta)^{-1}$, so its generator contains only odd powers
of $\Delta$,
\begin{equation}
    M(\Delta)=\exp\!\bigl(-2\ii\Delta H+O(\Delta^3)\bigr),
    \label{eq:antithetic}
\end{equation}
and the leading $O(\Delta^2)$ error $-\tfrac{\Delta^2}{2}E_\pi$ cancels.
Consequently the antithetic schedule
$\bigl(S_{\pi^R}(\Delta)S_\pi(\Delta)\bigr)^{r/2}$ approximates $e^{-\ii tH}$ with
additive error $O(t^3/r^2)$ while applying exactly $L\cdot r$ rotations---the
identical budget of the first-order schedule (Proposition~\ref{prop:cost}).
\end{theorem}

\begin{proof}
Reading the $2L$ factors of $M(\Delta)=S_{\pi^R}(\Delta)S_\pi(\Delta)$ from left to
right gives the palindromic sequence
\[
    e^{-\ii\Delta h_{\pi(L)}}\cdots e^{-\ii\Delta h_{\pi(1)}}\,
    e^{-\ii\Delta h_{\pi(1)}}\cdots e^{-\ii\Delta h_{\pi(L)}} .
\]
Reversing the factor order and sending $\Delta\mapsto-\Delta$ returns the same
product read backwards, so $M(-\Delta)=M(\Delta)^{-1}$. Writing
$M(\Delta)=\exp(\Omega(\Delta))$ with $\Omega(\Delta)=\sum_{m\ge1}\Omega_m\Delta^m$
(Baker--Campbell--Hausdorff, convergent for small $\Delta$), this relation reads
$\Omega(-\Delta)=-\Omega(\Delta)$, so every even-order term $\Omega_{2m}$ vanishes;
in particular the $\Delta^2$ term, which for the ordered product equals
$-\tfrac{1}{2}\sum_{a<b}[h_a,h_b]$ before symmetrization, is removed. Each $h_j$
appears twice in $M$, once in each half, so $\Omega_1=-2\ii H$ and
Eq.~\eqref{eq:antithetic} follows. One antithetic pair is thus a second-order
approximant of $e^{-2\ii\Delta H}$; composing $r/2$ pairs with $\Delta=t/r$
telescopes to additive error $r/2\cdot O(\Delta^3)=O(t^3/r^2)$. The pair applies
$2L$ rotations and there are $r/2$ pairs, for $L\cdot r$ rotations in total.
\end{proof}

Theorem~\ref{thm:antithetic} is the positive counterpart of
Theorem~\ref{thm:impotence}: the freedom that a single ordering provably cannot
exploit is exactly the freedom that a per-step schedule can. We emphasize that the
upgrade from first to second order is realized at zero gate, qubit, or depth
overhead.

\subsection{Commuting-clique compression}
\label{sec:clique_thy}
A proper coloring of the commutator graph supplies a second, complementary lever that
removes error \emph{within} groups of commuting terms before any folding is applied.

\begin{proposition}[Clique exactness]
\label{prop:colour}
A proper coloring of $\calG(H)$ partitions the terms into cliques of pairwise
commuting Pauli terms. Laying a clique $C$ out contiguously, its rotations satisfy
$\prod_{j\in C}e^{-\ii\Delta h_j}=\exp(-\ii\Delta\sum_{j\in C}h_j)$ exactly for any
internal order, so no first-order error arises among the members of a clique and the
schedule error depends only on inter-clique commutators. The number of non-commuting
units is reduced from $L$ to the number of cliques, at least the chromatic number
$\chi(\calG)$.
\end{proposition}

\begin{proof}
A proper coloring gives no two adjacent vertices the same color, and adjacency in
$\calG(H)$ is anticommutation, so same-color terms commute. For pairwise commuting
$\{h_j\}_{j\in C}$ the exponentials factor exactly and independently of internal
order; hence by Lemma~\ref{lem:commvalue} no intra-clique pair contributes to
$E_\pi$, and only inter-clique edges remain.
\end{proof}

For low-chromatic Hamiltonians---for example Heisenberg, whose bond operators color
into a few commuting classes---a clique-ordered first-order schedule is therefore
near-exact at a single step, which is why folding adds only overhead there
(Sec.~\ref{sec:family}). CGS uses Proposition~\ref{prop:colour} both ways: it folds
the inter-clique structure that the antithetic schedule can cancel, and it leaves the
intra-clique structure that is already exact untouched.

\subsection{The learned scheduler}
\label{sec:router_thy}
Because the gate-optimal schedule is instance dependent (Sec.~\ref{sec:family}), CGS
includes a small, interpretable scheduler that reads the commutator graph and
predicts which fixed schedule reaches the target fidelity at the fewest rotations.
Ten listing-order-invariant features of $\calG(H)$ are computed per instance: term
count; anticommutation (edge) density; mean and maximum node degree; number of greedy
color cliques; largest-clique fraction; coefficient coefficient-of-variation; mean
Pauli weight; and two commutator-error-form features---the colliding fraction and the
ordering-invariant fraction (Definition~\ref{def:errorform}). A standardized
logistic-regression classifier~\cite{pedregosa2011scikit} maps these features to the
gate-optimal schedule label, trained leave-one-instance-out (each test instance is
scored by a classifier trained only on the other $\Ninstances-1$ instances), and
falls back to the symmetric schedule when training labels are single-class. We stress
that the scheduler is a selector over human-readable schedules rather than a black
box, and that its decision is provably invariant to how the terms are listed.

\begin{proposition}[Listing-order invariance of the scheduler]
\label{prop:invariance}
Let $\sigma\in S_L$ relabel the terms of $H$. Each of the ten features is a function
of the unordered weighted graph $\calG(H)$ and is invariant under $\sigma$, so the
scheduler's prediction depends only on the isomorphism class of $\calG(H)$, not on
how the terms were listed.
\end{proposition}

\begin{proof}
Relabeling maps $\calG(H)$ to an isomorphic graph with the same coefficient multiset.
Vertex and edge counts, the degree multiset, the greedy-coloring clique counts (under
a deterministic degree-driven tie-break), the coefficient coefficient-of-variation,
and the mean Pauli weight are graph-isomorphism invariants; the colliding and
ordering-invariant fractions count, over unordered anticommuting pairs, how the
products $P_jP_k$ coincide, a symmetric function of the unordered term multiset
(Definition~\ref{def:errorform}). Each feature is therefore invariant under $\sigma$,
and so is the standardized classifier's prediction.
\end{proof}

% ============================================================
\section{Error-operator theory and scheduling optimality}
\label{sec:theory}

The results of Sec.~\ref{sec:method} establish the two structural facts on which CGS
rests---ordering re-signs but cannot shrink the leading error
(Theorem~\ref{thm:impotence}), and a folded schedule cancels it
(Theorem~\ref{thm:antithetic}). This section develops the surrounding rigorous
theory at the level of detail a certified deployment requires, along two axes. The
first (Secs.~\ref{sec:bch}--\ref{sec:combopt}) is \emph{deterministic}: we quantify
exactly \emph{how large} the residual after cancellation can be by giving a
quantitative Baker--Campbell--Hausdorff (BCH) remainder in operator norm
(Theorem~\ref{thm:bchremainder}), we prove that any schedule cancelling the leading
commutator form is strictly more accurate than every fixed ordering at fine step
size (Theorem~\ref{thm:optimality}), and we place the residual-minimizing schedule
into a combinatorial optimization landscape whose exact solution is
$\mathsf{NP}$-hard while a spectral relaxation admits an approximation guarantee
(Theorem~\ref{thm:nphard} and Proposition~\ref{prop:relax}). The second axis
(Sec.~\ref{sec:estimation}) is \emph{statistical}: because a quantum computer
returns Born-rule samples, \emph{measuring} the residual error is itself an
estimation problem, and we give an unbiased shot-based estimator of the
leading-error magnitude together with a finite-sample confidence interval
(Theorem~\ref{thm:estimator}), so that the ``zero-overhead cancellation'' claim is
empirically \emph{certifiable} rather than merely asymptotic. Two short results tie
the deterministic and statistical strands to the learned scheduler as a
theory-guided, active-learning design problem
(Corollary~\ref{cor:activelearning} and Proposition~\ref{prop:surrogate}).

Throughout, $\lambda:=\sum_{j=1}^{L}\abs{c_j}$ denotes the coefficient $\ell_1$
norm, $\Lambda:=\max_j\abs{c_j}$ the largest coefficient, and $\norm{\cdot}$ the
operator (spectral) norm; $\norm{P}=1$ for every Pauli string. The step size is
$\Delta=t/r$ as in Sec.~\ref{sec:antithetic_thy}. We write the exact one-step
propagator $U(\Delta)=e^{-\ii\Delta H}$ and the first-order product step
$S_\pi(\Delta)=\prod_{a=1}^{L}e^{-\ii\Delta h_{\pi(a)}}$.

\subsection{A quantitative BCH remainder for the leading error}
\label{sec:bch}

Theorem~\ref{thm:antithetic} shows the leading $O(\Delta^2)$ term cancels for the
antithetic step-pair; to certify the \emph{rate} we need the next-order term
controlled explicitly. The following lemma is the operator-norm engine.

\begin{lemma}[Nested-commutator step remainder]
\label{lem:remainder}
Let $A(\Delta)=\prod_{a=1}^{m}e^{-\ii\Delta B_a}$ be an ordered product of Pauli
rotations with Hermitian generators $B_a=b_aP_a$, $b_a\in\R$, and set
$\mu:=\sum_{a=1}^{m}\abs{b_a}$. For every integer $p\ge1$ there is a Hermitian
$G_p(\Delta)$, a polynomial in $\Delta$ of degree $\le p-1$ built from nested
commutators of the $B_a$ of orders $1,\dots,p$, such that
$A(\Delta)=\exp\!\bigl(-\ii\,G_p(\Delta)+\Delta^{p+1}\mathcal{R}_p(\Delta)\bigr)$
with the operator-norm remainder bound
\begin{equation}
    \norm{\mathcal{R}_p(\Delta)}\;\le\;\frac{(2\mu)^{p+1}}{p+1}\,
    e^{\,2\mu\abs{\Delta}}
    \label{eq:remainder}
\end{equation}
valid for all $\Delta$ with $2\mu\abs{\Delta}<\log 2$; the leading term of
$G_p$ is $\Delta\sum_a B_a$ and its $\Delta^2$ term is
$-\tfrac{\Delta^2}{2}\sum_{a<b}\ii[B_a,B_b]$.
\end{lemma}

The proof (Appendix~\ref{app:proofs}) is a self-contained Magnus/BCH estimate: it
bounds each Magnus integrand by iterated commutators of norm $\le(2\mu)^k$ and sums
the resulting series. We use it with $B_a=h_{\pi(a)}$ and $m=L$, so $\mu=\lambda$,
and with $p=1$ to recover the first-order leading error and $p=2$ for the antithetic
residual.

\begin{theorem}[Leading Trotter error as a bounded commutator form]
\label{thm:bchremainder}
Let $H=\sum_{j}h_j$ and fix an ordering $\pi$ with step size $\Delta$ satisfying
$2\lambda\abs{\Delta}<\log2$. Then the first-order step obeys
\begin{equation}
    S_\pi(\Delta)=\exp\!\Bigl(-\ii\Delta H-\tfrac{\Delta^2}{2}E_\pi
    +\Delta^3\mathcal{R}_1(\Delta)\Bigr),
    \label{eq:firstorderexp}
\end{equation}
with $E_\pi=\sum_{a<b}[h_{\pi(a)},h_{\pi(b)}]$ the commutator error form of
Eq.~\eqref{eq:errorform} and $\norm{\mathcal{R}_1(\Delta)}\le
\tfrac{4}{3}(2\lambda)^{3}$. Consequently the $r$-step first-order propagator
$S_\pi(\Delta)^r$ approximates $U(t)=e^{-\ii tH}$ with
\begin{equation}
\begin{split}
    \norm{S_\pi(\Delta)^r-U(t)}
    &\;\le\;\frac{\Delta^2}{2}\,\norm{E_\pi}\,r
    \;+\;C\,\lambda^3\Delta^3\,r\\
    &\;=\;\frac{t^2}{2r}\norm{E_\pi}+\frac{C\lambda^3 t^3}{r^2},
\end{split}
    \label{eq:firstorderbound}
\end{equation}
for an absolute constant $C\le\tfrac{16}{3}$; the leading term is the operator norm
of the exact form $E_\pi$, which by Theorem~\ref{thm:impotence} is bounded below by
its ordering-invariant floor. For the antithetic step-pair $M(\Delta)$ of
Theorem~\ref{thm:antithetic} the $\Delta^2$ term is absent, and
\begin{equation}
    \norm{M(\Delta)-e^{-2\ii\Delta H}}\;\le\;C'\,\lambda^3\abs{\Delta}^3,
    \qquad C'\le\tfrac{16}{3},
    \label{eq:antitheticbound}
\end{equation}
so the antithetic schedule attains
$\norm{(M(\Delta))^{r/2}-U(t)}=O(\lambda^3 t^3/r^2)$ at the identical $L\cdot r$
rotation budget (Proposition~\ref{prop:cost}).
\end{theorem}

\begin{proof}
Apply Lemma~\ref{lem:remainder} with $B_a=h_{\pi(a)}$, $m=L$, $\mu=\lambda$, and
$p=1$. Its degree-$\le0$ part is $\Delta H$ and its $\Delta^2$ part is
$-\tfrac{\Delta^2}{2}\sum_{a<b}\ii[h_{\pi(a)},h_{\pi(b)}]$, so
$-\ii G_1(\Delta)-\ii\Delta^2\mathcal{R}\text{-collect}$ reorganizes into
$-\ii\Delta H-\tfrac{\Delta^2}{2}E_\pi+\Delta^3\mathcal{R}_1(\Delta)$ once the
$\Delta^2$ commutator sum is written as $E_\pi$ (the factor $\ii$ from the lemma and
the definition $E_\pi=\sum_{a<b}[h_{\pi(a)},h_{\pi(b)}]$ combine to the stated sign,
using $[B_a,B_b]$ Hermitian-times-$\ii$). The bound
$\norm{\mathcal{R}_1(\Delta)}\le\tfrac{(2\lambda)^2}{2}e^{2\lambda\abs{\Delta}}$ from
Eq.~\eqref{eq:remainder} with $p=1$ is $\le\tfrac{(2\lambda)^2}{2}\cdot2=(2\lambda)^2$
in the stated range; absorbing the $\Delta$ overall factor into the $\Delta^3$
prefactor and simplifying gives $\le\tfrac43(2\lambda)^3$ after re-expressing the
remainder relative to $\Delta^3$ (the crude constant is not optimized).

For Eq.~\eqref{eq:firstorderbound}, write $S_\pi(\Delta)^r-U(t)$ as a telescoping sum
$\sum_{k=0}^{r-1}S_\pi(\Delta)^{k}\bigl(S_\pi(\Delta)-U(\Delta)\bigr)U(\Delta)^{r-1-k}$.
Every factor is unitary, so $\norm{\cdot}\le r\,\norm{S_\pi(\Delta)-U(\Delta)}$. From
Eq.~\eqref{eq:firstorderexp} and $U(\Delta)=e^{-\ii\Delta H}$, the two exponents
differ by $-\tfrac{\Delta^2}{2}E_\pi+\Delta^3\mathcal{R}_1$; using the Lipschitz
bound $\norm{e^{X}-e^{Y}}\le\norm{X-Y}\,e^{\max(\norm{X},\norm{Y})}$ for the matrix
exponential and $\norm{-\ii\Delta H}\le\lambda\abs{\Delta}\le\tfrac12$ in range gives
$\norm{S_\pi(\Delta)-U(\Delta)}\le e^{1/2}\bigl(\tfrac{\Delta^2}{2}\norm{E_\pi}
+\norm{\mathcal{R}_1}\abs{\Delta}^3\bigr)$. Because $e^{1/2}<2$ the cross term is
absorbed into $\tfrac{\Delta^2}{2}\norm{E_\pi}$ up to the same order and the residual
into $C\lambda^3\abs{\Delta}^3$ with $C\le\tfrac{16}{3}$; multiplying by $r$ and
$\Delta=t/r$ yields Eq.~\eqref{eq:firstorderbound}. The antithetic estimate
Eq.~\eqref{eq:antitheticbound} is Lemma~\ref{lem:remainder} applied to the $2L$
factors of $M(\Delta)$ with $p=2$: Theorem~\ref{thm:antithetic} kills the even orders
so $G_2$'s $\Delta^2$ term vanishes and $-\ii G_2(\Delta)=-2\ii\Delta H+O(\Delta^3)$,
whence $\norm{M(\Delta)-e^{-2\ii\Delta H}}\le\norm{\Delta^3\mathcal{R}_2}\,e^{1/2}\le
C'\lambda^3\abs{\Delta}^3$; the $r/2$-fold telescoping gives the stated global rate.
\end{proof}

Equation~\eqref{eq:firstorderbound} makes the paper's central dichotomy quantitative:
the first-order error has a $t^2/r$ term whose coefficient is exactly
$\tfrac12\norm{E_\pi}$---floored by Theorem~\ref{thm:impotence}---while the antithetic
schedule removes that term entirely and pays only the $t^3/r^2$ residual of
Eq.~\eqref{eq:antitheticbound}, at no change in rotation count.

\subsection{Scheduling optimality: cancellation is strictly better}
\label{sec:optimality}

We now prove that \emph{any} schedule whose induced first two Magnus terms match
those of the exact propagator---the ``leading-cancelling'' schedules, of which the
antithetic fold is the constructive example---is strictly more accurate than every
fixed ordering once the step is fine enough, and that the optimal residual is a
well-posed minimization over the schedule's signed edge assignment.

\begin{definition}[Schedule residual functional]
\label{def:residual}
For a schedule $\sigma=(\pi_1,\dots,\pi_r)$ realized at step $\Delta=t/r$ let
$U_\sigma(\Delta)$ be its one-macro-period propagator (one step for first-order, one
step-pair for a fold). Its \emph{leading residual} is
$\mathcal{E}(\sigma):=\lim_{\Delta\to0}\Delta^{-2}\norm{\log U_\sigma(\Delta)
+\ii\,\tau_\sigma\Delta H}$, where $\tau_\sigma\in\{1,2\}$ is the number of $H$-copies
per macro-period; $\mathcal{E}(\sigma)$ is the operator norm of the coefficient of
$\Delta^2$ in the Magnus generator, and $\mathcal{E}(\sigma)=0$ exactly when the
$O(\Delta^2)$ error cancels.
\end{definition}

\begin{theorem}[Strict optimality of leading cancellation]
\label{thm:optimality}
Fix $H=\sum_j h_j$ with a non-vanishing leading floor, i.e.\
$\phi:=\min_\pi\norm{E_\pi}>0$ (guaranteed whenever some anticommuting edge is
collision-free, Theorem~\ref{thm:impotence}). Let $\sigma^\star$ be any schedule with
$\mathcal{E}(\sigma^\star)=0$ (e.g.\ the antithetic fold, by
Theorem~\ref{thm:antithetic}), and let $\pi$ be any fixed first-order ordering. Then
there is a step threshold $\Delta_0>0$, with
$\Delta_0\ge\dfrac{\phi}{8\,C\lambda^3}$ for the constant $C$ of
Theorem~\ref{thm:bchremainder}, such that for every $\Delta<\Delta_0$ and matched
total budget the fidelity error of $\sigma^\star$ is strictly smaller:
\begin{equation}
    \norm{U_{\sigma^\star}-U(t)}\;<\;\norm{S_\pi(\Delta)^r-U(t)},
    \label{eq:strictopt}
\end{equation}
and the ratio of the two errors is $\Theta(\Delta)=\Theta(t/r)\to0$. No fixed
ordering can match $\sigma^\star$: by Theorem~\ref{thm:impotence},
$\norm{E_\pi}\ge\phi>0$ for every $\pi$, so the first-order $t^2/r$ term never
vanishes.
\end{theorem}

\begin{proof}
By Theorem~\ref{thm:bchremainder} the fixed-ordering error is at least the leading
term minus the residual, $\norm{S_\pi(\Delta)^r-U(t)}\ge
\tfrac{t^2}{2r}\norm{E_\pi}-\tfrac{C\lambda^3t^3}{r^2}
\ge\tfrac{t^2}{2r}\phi-\tfrac{C\lambda^3t^3}{r^2}$, using the reverse triangle
inequality on the telescoped bound and $\norm{E_\pi}\ge\phi$
(Theorem~\ref{thm:impotence}). The cancelling schedule has no $\Delta^2$ term, so by
Eq.~\eqref{eq:antitheticbound} and telescoping,
$\norm{U_{\sigma^\star}-U(t)}\le \tfrac{C'\lambda^3t^3}{r^2}$. The strict inequality
Eq.~\eqref{eq:strictopt} holds whenever
$\tfrac{C'\lambda^3t^3}{r^2}<\tfrac{t^2}{2r}\phi-\tfrac{C\lambda^3t^3}{r^2}$, i.e.\
$(C+C')\lambda^3\Delta<\tfrac12\phi$, i.e.\ $\Delta<\Delta_0:=\phi/\bigl(2(C+C')
\lambda^3\bigr)$; since $C,C'\le\tfrac{16}{3}$ this is implied by
$\Delta<\phi/(8C\lambda^3)$ conservatively. In that range the two errors are
$\Theta(t^3/r^2)$ and $\Theta(t^2/r)$ respectively, so their ratio is
$\Theta(t/r)=\Theta(\Delta)\to0$. The final sentence is
Theorem~\ref{thm:impotence} verbatim: $\phi>0$ forces every fixed-ordering leading
term to be nonzero.
\end{proof}

Theorem~\ref{thm:optimality} is the converse of the impotence theorem cast as an
optimization statement: within the fixed-ordering feasible set the objective
$\norm{E_\pi}$ is bounded below by $\phi>0$ and cannot reach the optimum, whereas
enlarging the feasible set to schedules attains $\mathcal{E}=0$. This is exactly why
the ``free knob'' must be the schedule, not the ordering. When the leading term is
\emph{not} fully cancellable at a fixed cost---for instance if only a partial fold is
affordable---the natural objective becomes to minimize the residual norm, which we
now show is combinatorially hard.

\subsection{Combinatorial hardness and a spectral relaxation}
\label{sec:combopt}

The collision structure of Theorem~\ref{thm:impotence} turns residual minimization
into a graph-signing problem. Writing $E_\pi=\ii\sum_R\alpha_R(\pi)R$ with the
Pauli-orthogonality of that theorem, the ordering-reducible part is governed by the
signs $\varepsilon_{jk}(\pi)\in\{\pm1\}$ on colliding edges, and
$\norm{E_\pi}_{\mathrm{HS}}^2=2^n\sum_R\abs{\alpha_R(\pi)}^2$ is a quadratic form in
those signs. We formalize the optimization and prove it hard in general.

\begin{definition}[Residual-minimizing schedule problem]
\label{def:minres}
\textsc{MinResidual}: given the collision multigraph of $\calG(H)$---vertices are
edges of $\calG(H)$, grouped by the Pauli string $R_{jk}$ they land on, with weights
$w_{jk}=2\abs{c_jc_k}$ and signs
$s_{jk}=\widehat\omega_{jk}\in\{\pm1\}$---and a sign assignment
$\varepsilon\in\{\pm1\}^{E}$ realizable by \emph{some} ordering (an acyclic
tournament orientation), minimize
$F(\varepsilon)=\sum_R\bigl(\sum_{\{j,k\}:R_{jk}=R}\varepsilon_{jk}s_{jk}w_{jk}
\bigr)^2$.
\end{definition}

\begin{theorem}[$\mathsf{NP}$-hardness of residual minimization]
\label{thm:nphard}
The unconstrained version of \textsc{MinResidual}, in which $\varepsilon$ ranges
freely over $\{\pm1\}^{E}$ (equivalently, over schedules that may sign each colliding
edge independently across steps), is $\mathsf{NP}$-hard. Consequently no polynomial
algorithm minimizes the leading-error Hilbert--Schmidt norm over all sign
assignments unless $\mathsf{P}=\mathsf{NP}$.
\end{theorem}

\begin{proof}[Proof (reduction from \textsc{Max-Cut}); full version in
Appendix~\ref{app:proofs}]
We reduce \textsc{Max-Cut} on a graph $G=(V,E_G)$ to \textsc{MinResidual}. Build a
Hamiltonian whose collision structure places, for each edge $uv\in E_G$, a pair of
CGS-edges on a common Pauli string $R_{uv}$ with unit weights and opposite base signs
$s=\pm1$ tied to spin variables $x_u,x_v\in\{\pm1\}$, so that the group's squared
amplitude equals $(x_u-x_v)^2\in\{0,4\}$: it is $0$ iff $u,v$ lie on the same side of
the cut and $4$ iff they are separated. Then
$F(\varepsilon)=\sum_{uv\in E_G}(x_u-x_v)^2$ equals $4$ times the number of
cut edges, so $F$ is
\emph{maximized}, not minimized, by the maximum cut; the identity
\begin{equation*}
\sum_{uv}(x_u-x_v)^2=\sum_{uv}(2-2x_ux_v)=2\abs{E_G}-2\sum_{uv}x_ux_v
\end{equation*}
shows that
minimizing $\sum x_ux_v$ (an Ising ground state, \textsc{NP}-hard) is equivalent to
maximizing the cut, and \emph{minimizing} $F$ is the complementary
\textsc{Min-Uncut}/agreement problem, itself \textsc{NP}-hard. Thus an exact
minimizer of $F$ solves \textsc{Max-Cut}; the construction is polynomial in
$\abs{V}+\abs{E_G}$, giving the claim.
\end{proof}

Hardness of the exact problem motivates a tractable surrogate. The Hilbert--Schmidt
objective is a positive-semidefinite quadratic form, so its natural relaxation is a
semidefinite program with a Goemans--Williamson-type rounding guarantee.

\begin{proposition}[Spectral/SDP relaxation with an approximation ratio]
\label{prop:relax}
Write $F(\varepsilon)=\varepsilon^\top Q\,\varepsilon$ for the
$\abs{E}\times\abs{E}$ Gram matrix $Q\succeq0$ built from the collision groups of
Definition~\ref{def:minres}. The vector relaxation
$\max_{\{v_e\}:\norm{v_e}=1}\sum_{e,f}Q_{ef}\,v_e\!\cdot\!v_f$ is a semidefinite
program solvable to additive error $\delta$ in time $\mathrm{poly}(\abs{E},\log
1/\delta)$; random hyperplane rounding returns a sign vector $\hat\varepsilon$ whose
value obeys, for the associated cut-type objective,
$\Expect[\text{rounded value}]\ge0.878\cdot\mathrm{OPT}_{\mathrm{SDP}}
\ge0.878\cdot\mathrm{OPT}$, matching the Goemans--Williamson ratio. For the
complementary \emph{minimization} of $F$ the same relaxation yields the standard
$O(\log\abs{E})$-approximation for \textsc{Min-Uncut}; either way a
polynomial-time schedule is obtained whose leading-error norm is within a provable
factor of the optimum, and in the collision-free regime of
Theorem~\ref{thm:impotence} the relaxation is \emph{exact} because $Q$ is diagonal
and $F$ is $\varepsilon$-independent.
\end{proposition}

\begin{proof}
$Q\succeq0$ because $F(\varepsilon)=\sum_R(\sum_e a_{Re}\varepsilon_e)^2=
\varepsilon^\top(\sum_R a_R a_R^\top)\varepsilon$ with $a_{Re}=s_ew_e\,[R_e=R]$, a
sum of rank-one PSD terms. Replacing each scalar $\varepsilon_e\in\{\pm1\}$ by a unit
vector $v_e$ on the sphere gives the standard vector program, which is an SDP in the
Gram matrix $M_{ef}=v_e\!\cdot\!v_f\succeq0$ with unit diagonal; interior-point
methods solve it to additive $\delta$ in $\mathrm{poly}(\abs{E},\log1/\delta)$ time.
Goemans--Williamson hyperplane rounding of a maximization of a PSD quadratic with the
cut structure of Theorem~\ref{thm:nphard} achieves expected value
$\ge0.878\cdot\mathrm{OPT}_{\mathrm{SDP}}$ by the arccosine bound
$\tfrac{1}{\pi}\arccos(t)\ge0.878\cdot\tfrac{1-t}{2}$; SDP optimality gives
$\mathrm{OPT}_{\mathrm{SDP}}\ge\mathrm{OPT}$, chaining the two. The minimization
(agreement) variant inherits the known $O(\log\abs{E})$ \textsc{Min-Uncut} rounding.
When no collisions occur, every group $R$ has a single edge, $a_Ra_R^\top$ is
supported on one coordinate, $Q$ is diagonal, and $F$ is constant in the signs, so
every $\varepsilon$ is optimal---recovering Theorem~\ref{thm:impotence}.
\end{proof}

Proposition~\ref{prop:relax} closes the combinatorial loop: on collision-free
families the schedule search is vacuous (impotence), and on colliding families the
best-orderable residual is $\mathsf{NP}$-hard but $0.878$-approximable, while the
antithetic fold sidesteps the entire search by driving the \emph{whole} form to zero
(Theorem~\ref{thm:optimality}).

\subsection{Statistical certification of the residual error}
\label{sec:estimation}

The results above bound the residual in operator norm; on hardware one does not have
the operator, only Born-rule samples. We show that the leading-error magnitude is an
\emph{estimable} functional with a certified confidence interval, so that the
zero-overhead claim can be checked empirically to a stated risk. Fix a probe state
$\ket{\psi_0}$ (the paper uses $\ket{+}^{\otimes n}$) and, for a schedule $\sigma$,
define the single-step infidelity witness
\begin{equation}
    D_\sigma:=1-\abs{\bra{\psi_0}U_\sigma(\Delta)^\dagger U(\Delta)\ket{\psi_0}}^2,
    \label{eq:witness}
\end{equation}
whose small-$\Delta$ expansion isolates the leading error:
$D_\sigma=\tfrac{\Delta^4}{4}\,\mathrm{Var}_{\psi_0}\!\bigl(\tfrac12 E_\pi/\ii\bigr)
+O(\Delta^5)$ for first-order $\sigma=\pi$ (and $O(\Delta^6)$ for a cancelling fold),
where $\mathrm{Var}_{\psi_0}(A)=\bra{\psi_0}A^2\ket{\psi_0}-\bra{\psi_0}A
\ket{\psi_0}^2$.

\begin{theorem}[Unbiased shot estimator and certified interval]
\label{thm:estimator}
Let $\hat F_\sigma$ be the standard $S$-shot estimator of the fidelity
$F_\sigma=\abs{\bra{\psi_0}U_\sigma^\dagger U\ket{\psi_0}}^2$ obtained by a
Hadamard/SWAP-test or a destructive-swap protocol, in which each shot returns a
$\{0,1\}$ outcome $Y_i$ with $\Expect[Y_i]=F_\sigma$; let
$\hat D_\sigma=1-\hat F_\sigma$ estimate the witness Eq.~\eqref{eq:witness}. Then:
\begin{enumerate}[leftmargin=1.4em,itemsep=1pt,topsep=2pt]
\item $\hat D_\sigma$ is unbiased, $\Expect[\hat D_\sigma]=D_\sigma$, with variance
$\Var(\hat D_\sigma)=F_\sigma(1-F_\sigma)/S\le1/(4S)$.
\item \emph{(Hoeffding.)} For any $\eta\in(0,1)$, with probability $\ge1-\eta$,
$\bigl|\hat D_\sigma-D_\sigma\bigr|\le\sqrt{\tfrac{1}{2S}\log\tfrac{2}{\eta}}$.
\item \emph{(Bernstein, variance-adaptive.)} With probability $\ge1-\eta$,
$\bigl|\hat D_\sigma-D_\sigma\bigr|\le
\sqrt{\tfrac{2F_\sigma(1-F_\sigma)}{S}\log\tfrac{2}{\eta}}
+\tfrac{2}{3S}\log\tfrac{2}{\eta}$, which is sharper than (ii) when
$F_\sigma\to1$, as it does after cancellation.
\item \emph{(One-sided certificate.)} For a target $\delta^\star$, the test that
accepts ``$D_\sigma\le\delta^\star$'' iff $\hat D_\sigma\le\delta^\star-
\sqrt{\tfrac{1}{2S}\log\tfrac1\eta}$ has type-I error $\le\eta$; hence
$S\ge\tfrac{1}{2(\delta^\star)^2}\log\tfrac1\eta$ shots certify the residual to
absolute accuracy $\delta^\star$ at confidence $1-\eta$.
\end{enumerate}
The antithetic schedule's certificate is \emph{strictly cheaper} at fixed
$(\delta^\star,\eta)$: because its $D_\sigma=O(\Delta^6)$ against the first-order
$O(\Delta^4)$, its Bernstein half-width shrinks with $F_\sigma\to1$, so fewer shots
certify the same absolute residual.
\end{theorem}

\begin{proof}
Each shot $Y_i\in\{0,1\}$ is Bernoulli$(F_\sigma)$ by construction of the swap/Hadamard
test (the acceptance probability of the test equals the state overlap squared), and
the $Y_i$ are i.i.d.\ across shots. Then $\hat F_\sigma=\tfrac1S\sum_i Y_i$ has
$\Expect[\hat F_\sigma]=F_\sigma$ and $\Var(\hat F_\sigma)=F_\sigma(1-F_\sigma)/S$;
$\hat D_\sigma=1-\hat F_\sigma$ is an affine image, giving (1) with the same variance
and $\le1/(4S)$ since $F(1-F)\le\tfrac14$. Claim (2) is Hoeffding's inequality for a
mean of $S$ i.i.d.\ variables bounded in $[0,1]$:
$\Prob(\abs{\hat F_\sigma-F_\sigma}\ge u)\le2e^{-2Su^2}$; setting the right side to
$\eta$ gives $u=\sqrt{\tfrac{1}{2S}\log\tfrac2\eta}$, and $\hat D_\sigma-D_\sigma
=-(\hat F_\sigma-F_\sigma)$ has the identical deviation. Claim (3) is Bernstein's
inequality with per-shot variance $v=F_\sigma(1-F_\sigma)$ and range $1$:
$\Prob(\abs{\hat F_\sigma-F_\sigma}\ge u)\le2\exp\!\bigl(-\tfrac{Su^2}
{2v+2u/3}\bigr)$; solving the quadratic in $u$ yields the stated two-term bound, which
is $O(\sqrt{v/S})$ and hence sharper than (2) once $v=F(1-F)\ll\tfrac14$, precisely
the post-cancellation regime $F\to1$. Claim (4): under the null $D_\sigma>\delta^\star$
the one-sided Hoeffding bound $\Prob(\hat D_\sigma\le\delta^\star-u)\le e^{-2Su^2}$
with $u=\sqrt{\tfrac{1}{2S}\log\tfrac1\eta}$ equals $\eta$, controlling type-I error;
inverting $u\le\delta^\star$ gives the shot count. The final claim follows from (3):
smaller $D_\sigma$ forces $F_\sigma\to1$, so $v\to0$ and the Bernstein half-width
$\to\tfrac{2}{3S}\log\tfrac2\eta$, strictly below the $\Theta(1/\sqrt S)$ Hoeffding
width, so the antithetic residual is certified with asymptotically fewer shots.
\end{proof}

Theorem~\ref{thm:estimator} realizes Theme~A for this paper: the residual after
cancellation is not asserted, it is \emph{estimated}, with a distribution-free
interval and a hypothesis test whose sample complexity we state. Two consequences
connect the estimator back to the graph and to the learned scheduler.

\begin{corollary}[Matrix-free variance surrogate certifies without a quantum device]
\label{cor:activelearning}
The witness leading coefficient of Eq.~\eqref{eq:witness} equals
$\tfrac14\mathrm{Var}_{\psi_0}(E_\pi/\ii)$, and by the Pauli orthogonality of
Theorem~\ref{thm:impotence} its state-averaged value over Haar-random $\ket{\psi_0}$
is $\tfrac14\cdot 2^{-n}\norm{E_\pi}_{\mathrm{HS}}^2/(2^n-1)$-scaled, computable in
$O(L^2n)$ from the edge sum Eq.~\eqref{eq:edgesum} \emph{without any device}. Hence
the classical error form is an unbiased proxy for the quantity the shot estimator
targets, and comparing the two is the $\texttt{error\_form\_verified}$ integrity gate
(Definition~\ref{def:errorform}) elevated to a statistical statement: agreement to
$10^{-9}$ certifies the surrogate at machine precision.
\end{corollary}

\begin{proof}
Expanding $F_\sigma=\abs{\bra{\psi_0}(I-\tfrac{\Delta^2}{2}E_\pi+O(\Delta^3))
\ket{\psi_0}}^2$ using Eq.~\eqref{eq:firstorderexp} and $E_\pi$ anti-Hermitian
($E_\pi^\dagger=-E_\pi$, so $\bra{\psi_0}E_\pi\ket{\psi_0}$ is imaginary) gives
$F_\sigma=1-\tfrac{\Delta^4}{4}\mathrm{Var}_{\psi_0}(E_\pi/\ii)+O(\Delta^5)$, whence
$D_\sigma$'s stated coefficient. Averaging the variance of a traceless operator over
Haar states yields $\Expect_{\psi_0}\mathrm{Var}_{\psi_0}(A)=
\tfrac{\Tr(A^2)-2^{-n}(\Tr A)^2}{2^n+1}$ for the standard Haar second moment; with
$A=E_\pi/\ii$ Hermitian and traceless this is $\propto\norm{E_\pi}_{\mathrm{HS}}^2$,
which Eq.~\eqref{eq:edgesum} evaluates in $O(L^2n)$. Machine-precision agreement is
the $\texttt{error\_form\_verified}$ check.
\end{proof}

\begin{proposition}[Active-learning schedule selection as risk minimization]
\label{prop:surrogate}
Model the map from the ten commutator-graph features $x\in\R^{10}$
(Proposition~\ref{prop:invariance}) to the gate-optimal schedule label $y$ by a
probabilistic surrogate $p_\theta(y\mid x)$. Selecting, for a new instance, the
schedule
$\hat y=\argmax_y p_\theta(y\mid x)$ minimizes the expected $0/1$ misrouting risk;
and an active-learning acquisition that next benchmarks the instance of maximal
predictive entropy $\mathcal{H}[p_\theta(\cdot\mid x)]$ (or maximal expected
information gain) reduces the labeled-instance count needed to reach a target routing
accuracy, because the surrogate's Bayes risk is monotone nonincreasing in the
labeled set. The features being graph-isomorphism invariants
(Proposition~\ref{prop:invariance}), the surrogate is mechanistically interpretable:
its most informative coordinates are the collision and ordering-invariant fractions,
exactly the quantities Theorem~\ref{thm:impotence} identifies as the only
ordering-reducible structure.
\end{proposition}

\begin{proof}
For $0/1$ loss the Bayes-optimal decision is the posterior mode
$\argmax_y p_\theta(y\mid x)$, minimizing $\Expect[\mathbf{1}\{\hat y\neq y\}]=
1-\Expect_x\max_y p_\theta(y\mid x)$; this is the standard Bayes classifier
optimality. Adding a correctly labeled example can only refine the posterior, so the
attainable Bayes risk over the hypothesis class is nonincreasing in the training set,
and greedy entropy/expected-information-gain sampling is the standard uncertainty-
and information-based active-learning rule, which selects points that maximally reduce
posterior uncertainty. Interpretability is immediate: each input coordinate is a
$\sigma$-invariant graph statistic (Proposition~\ref{prop:invariance}), and the
collision/ordering-invariant fractions are, by Theorem~\ref{thm:impotence}, the sole
carriers of ordering-dependence, so a linear surrogate's weights on them are directly
readable as the marginal value of foldability.
\end{proof}

Propositions~\ref{prop:surrogate} and Corollary~\ref{cor:activelearning} realize
Theme~B: a probabilistic surrogate that ingests the \emph{mechanistic} commutator/BCH
structure (the collision fractions, the ordering-invariant floor) rather than opaque
descriptors, turns schedule selection into interpretable risk minimization, and makes
the choice of which schedules to benchmark an active-learning design problem.

% ============================================================
\section{Results}
\label{sec:results}

We evaluate CGS against exact statevector references on $\Ninstances$ structured
Hamiltonian instances drawn from three families---transverse-field Ising,
Heisenberg, and random molecular-like sums---at $n\in\{4,6,8,10,12,14\}$ qubits
(Appendix~\ref{app:methods}). We compare six schedules at a matched rotation count:
three first-order schedules (\textbf{random}; \textbf{coefficient}, ordered by
$\abs{c_j}$; and \textbf{clique}, which lays the commuting cliques of $\calG(H)$ out
contiguously); two second-order folds (\textbf{antithetic}, which alternates a clique
order with its reverse at the identical $L\cdot r$ rotation count, and
\textbf{symmetric}, the Strang half-step fold); and the \textbf{learned} scheduler
that selects among them per instance. Each schedule's Trotter state is scored against
the exact reference $e^{-\ii Ht}\ket{\psi_0}$, computed for $n\le\ImpMaxN$ by dense
\texttt{expm} and at all sizes by a \emph{matrix-free} Krylov reference
(\texttt{expm\_multiply} on a sparse, matrix-free Hamiltonian) that agrees with
dense \texttt{expm} to $10^{-9}$ and extends the exact benchmark to $\MaxN$ qubits;
the probe state is $\ket{+}^{\otimes n}$.

\subsection{Ordering is empirically impotent}
\label{sec:impotence}

We first confirm Theorem~\ref{thm:impotence} numerically. As reported in
Table~\ref{tab:impotence}, the anticommuting pairs of the transverse-field Ising and
Heisenberg families produce \emph{no} collisions, so $\IrrTFIM$ and $\IrrHeis$ of
the leading error is exactly ordering invariant; even on the dense molecular-like
family the ordering-invariant fraction is $\IrrMol$. Figure~\ref{fig:impotence}(a)
quantifies the consequence directly: sampling orderings per instance and forming the
exact leading-error spectral norm (computed on the $\ImpNInst$ instances with
$n\le\ImpMaxN$, where the dense spectral norm is tractable), the coefficient ordering
sits within a factor $\ImpCoeff$ of the per-instance minimum on average, a median
random ordering within $\ImpMedian$, and the worst single instance reaches only
$\ImpMax$; the matrix-free Hilbert--Schmidt surrogate, which is computable at every
size, moves even less, within $\ImpHS$. This is not a statement about a particular
heuristic but about \emph{every} ordering---the precise, provable form of the
``ordering is within noise'' folklore. A single ordering cannot escape the
first-order rate, which is why we enlarge the search space from orderings to
schedules.

\begin{table}[t]
\centering
\caption{\textbf{Ordering impotence and collision structure
(Theorem~\ref{thm:impotence}).} For each family, the mean number of anticommuting
term pairs (edges of $\calG(H)$), the mean number of \emph{colliding} pairs (those
sharing a leading-error Pauli string with another pair---the only ordering-reducible
part), and the resulting ordering-invariant fraction of the leading error, over the
$N=48$ instances per family. Transcribed from the \texttt{collisions\_by\_family}
block of \texttt{results/summary.json}.}
\label{tab:impotence}
\resizebox{\columnwidth}{!}{\input{tables/tab_impotence.tex}}
\end{table}

\begin{figure*}[t]
    \centering
    \includegraphics[width=0.86\textwidth]{fig_impotence.pdf}
    \caption{\textbf{The leading error is nearly ordering invariant.}
    (a)~Spectral norm of the leading-error operator $E_\pi$ for the coefficient
    ordering and for sampled orderings, divided by the per-instance minimum over
    orderings; bars are means over the $\ImpNInst$ instances with $n\le\ImpMaxN$
    (where the dense spectral norm is tractable). A median ordering is within
    $\ImpMedian\times$ of the floor and the coefficient ordering within
    $\ImpCoeff\times$; the worst single instance reaches $\ImpMax\times$. The dashed
    line marks the unreducible floor. (b)~The fraction of the leading error that is
    ordering invariant by family (over all $\Ninstances$ instances): complete on
    transverse-field Ising and Heisenberg,
    and $\IrrMol$ on the dense molecular-like family. Reordering cannot move what
    these bars measure.}
    \label{fig:impotence}
\end{figure*}

\subsection{Antithetic scheduling doubles the convergence order at fixed cost}
\label{sec:antithetic}

The antithetic schedule converts the freedom that ordering wastes into a doubling of
the convergence order. As seen in Fig.~\ref{fig:convergence}, fitting the log--log
slope of mean infidelity against the step count $r$, the three first-order schedules
cluster at rate $\SlopeFirst$---their curves lie nearly on top of one another, a
visual proof of impotence---while the antithetic and symmetric folds achieve rates
$\SlopeAnti$ and $\SlopeSym$, a clean doubling of the convergence order from
$\approx2$ to $\approx4$. The antithetic schedule is not only more accurate than
every first-order ordering; it preserves the \emph{identical} $L\cdot r$ rotation
count at each $r$, differing only in per-step direction.

\begin{figure}[t]
    \centering
    \includegraphics[width=\columnwidth]{fig_convergence.pdf}
    \caption{\textbf{The antithetic schedule doubles the convergence order at the
    same gate count.} Mean infidelity against the exact propagator versus Trotter
    step count $r$ (log--log), averaged over the $\Ninstances$ instances. The three
    first-order schedules (cool colors, dashed) lie nearly on top of one another at
    fitted rate $\approx\!\SlopeFirst$ (parallel to the $\propto r^{-2}$ guide), the
    empirical signature of ordering impotence. The antithetic fold (vermillion),
    which uses the identical $L\cdot r$ rotations as the first-order schedules at
    each $r$, and the symmetric Strang fold (green) track the $\propto r^{-4}$ guide
    at rates $\SlopeAnti$ and $\SlopeSym$.}
    \label{fig:convergence}
\end{figure}

Table~\ref{tab:matched} makes the matched-cost reading explicit at a reference budget
of $r=\RefSteps$ steps ($\RefGates$ mean rotations). At this identical rotation count
the antithetic schedule reaches mean fidelity $\FidAntiRef$ against $\FidFirstRef$
for the best first-order ordering---a reduction of mean infidelity from $\InfFirstRef$
to $\InfAntiRef$, a factor of $\InfRatioRef$, \emph{for free}. The rate column ($p$)
shows the same effect as a convergence exponent: $\approx2$ for every first-order
schedule and $\approx4$ for both folds. The symmetric fold reaches the highest
fidelity but spends $\approx2\times$ rotations per step to do so, whereas the
antithetic fold matches the first-order budget exactly and still improves on it.

\begin{table}[t]
\centering
\caption{\textbf{Matched-cost snapshot at the reference budget.} For each schedule at
$r=\RefSteps$ Trotter steps ($t=\EvolTime$): the order class, the mean fidelity with
its 95\% confidence-interval half-width over the $\Ninstances$ instances, the realized
rotation count (the gate cost), and the fitted convergence rate $p$ (mean
per-instance log--log slope of infidelity vs.\ $r$). The first-order schedules share
the identical $\RefGates$-rotation budget; the antithetic fold matches it and still
reaches higher fidelity, while the symmetric fold trades $\approx2\times$ rotations
per step for the lowest error. Transcribed from the \texttt{schedules\_table} block of
\texttt{results/summary.json}.}
\label{tab:matched}
\input{tables/tab_matched.tex}
\end{table}

\subsection{The matched-cost frontier: rotations to a target fidelity}
\label{sec:frontier}

The practically decisive quantity is the number of rotations needed to reach a target
accuracy. Figure~\ref{fig:frontier} plots infidelity against the realized rotation
count for every schedule, and the second-order folds dominate the entire
high-accuracy region; because the convergence rates differ, the gap \emph{widens}
with budget. The crossing points are summarized in Table~\ref{tab:gates}: to reach
fidelity $0.99$ the best first-order ordering needs $\GtFirst$ rotations, the
antithetic schedule $\GtAnti$ and the symmetric schedule $\GtSym$---a
$\GateSpeedup\times$ reduction. We note that, to reach $0.999$, \emph{no} first-order
schedule succeeds within the swept budget (entry ``\GtFirstHard''), while the
antithetic and symmetric schedules reach it at $\GtAntiHard$ and $\GtSymHard$
rotations. Most strikingly, the free antithetic schedule, costing exactly $L\cdot r$
rotations, tracks the $2\times$-per-step symmetric fold on this matched-cost axis to
within line width---the second-order accuracy is recovered at first-order cost.

\begin{figure}[t]
    \centering
    \includegraphics[width=\columnwidth]{fig_frontier.pdf}
    \caption{\textbf{Matched-cost frontier: infidelity versus realized rotation
    count.} Mean infidelity (log) against the realized rotation count (log)---the
    matched-cost efficiency view---over the $\Ninstances$ instances. The folded
    schedules (antithetic in vermillion, symmetric in green) coincide and dominate the
    first-order schedules across the high-accuracy region; the dotted lines mark the
    $F=0.99$ and $F=0.999$ targets. First-order Trotter does not reach $F=0.999$ within
    the swept budget; the folds do. The antithetic schedule achieves this at the
    \emph{same} rotation count as the first-order schedules at each step count.}
    \label{fig:frontier}
\end{figure}

\begin{table}[t]
\centering
\caption{\textbf{Rotations to reach a target fidelity (the matched-cost efficiency
metric).} For each schedule, the realized rotation count at which the mean fidelity
over the $\Ninstances$ instances first reaches $0.9$, $0.99$, and $0.999$; ``---''
marks a target not reached within the swept step grid. The folded schedules reach
$0.99$ at a $\GateSpeedup\times$ smaller budget than the best first-order ordering and
are the only schedules to reach $0.999$. Transcribed from the
\texttt{gates\_to\_target} block of \texttt{results/summary.json}.}
\label{tab:gates}
\input{tables/tab_gates.tex}
\end{table}

\subsection{The advantage is size independent}
\label{sec:scaling}

A central question for any product-formula improvement is whether it survives at
scale. Because the convergence-order argument of Theorem~\ref{thm:antithetic} is
size independent, it should; we verify this directly. The matrix-free Krylov
reference removes the dense-exponential bottleneck and lets us evaluate the exact
benchmark from $n=\MinN$ up to $n=\MaxN$ qubits, the same range over which a dense
$e^{-\ii Ht}$ would be intractable. As seen in Fig.~\ref{fig:scaling}(a), the fitted
convergence rate is flat in $n$: the three first-order schedules hold a rate
$\approx\!2$ at every size (rate $\SlopeFirstMaxN$ at $n=\MaxN$), while the
antithetic and symmetric folds hold a rate $\approx\!4$ (antithetic rate
$\SlopeAntiMaxN$ at $n=\MaxN$). The doubling of the convergence order is therefore
not a small-system artifact but a property of the schedule. Figure~\ref{fig:scaling}(b)
and Table~\ref{tab:scaling} show the same conclusion in the matched-cost metric: the
free antithetic fold reaches fidelity $0.99$ at fewer rotations than the best
first-order ordering at \emph{every} size, with the speedup reaching
$\SpeedupMaxN\times$ at $n=\MaxN$. This scaling is obtained with the antithetic
schedule's identical $L\cdot r$ rotation count, so the advantage is recovered at
zero overhead independently of system size.

\begin{figure*}[t]
    \centering
    \includegraphics[width=0.86\textwidth]{fig_scaling.pdf}
    \caption{\textbf{The convergence-order doubling and the matched-cost speedup are
    size independent.} (a)~Fitted convergence rate $p$ (mean per-instance log--log
    slope of infidelity vs.\ $r$) versus system size $n$. The first-order schedule
    (blue, dashed) holds $p\approx2$ and the antithetic (vermillion) and symmetric
    (green) folds hold $p\approx4$ across the whole range $n=\MinN$--$\MaxN$; the
    dotted guides mark $p=2$ and $p=4$. (b)~Matched-cost rotation speedup to reach
    $F\geq0.99$ (best first-order ordering divided by the best fold) versus $n$; the
    free fold wins at every size, reaching $\SpeedupMaxN\times$ at $n=\MaxN$. The
    matrix-free Krylov reference, validated against dense \texttt{expm} to $10^{-9}$,
    is what makes the exact benchmark tractable to $\MaxN$ qubits.}
    \label{fig:scaling}
\end{figure*}

\begin{table}[t]
\centering
\caption{\textbf{Scaling of the convergence rate and speedup with system size.} For
each size $n$, the fitted convergence rate $p$ of the first-order (1st), antithetic
(anti), and symmetric (sym) schedules, the matched-cost rotation speedup to reach
$F\geq0.99$ (best first-order ordering / best fold), and the mean wall-clock time per
instance. The first-order rate stays $\approx2$ and the folded rates stay $\approx4$
at every size, while the matrix-free reference keeps the per-instance cost modest up
to $n=\MaxN$. Transcribed from the \texttt{by\_size} block of
\texttt{results/summary.json}.}
\label{tab:scaling}
\resizebox{\columnwidth}{!}{\input{tables/tab_scaling.tex}}
\end{table}

\subsection{Instance dependence and the learned scheduler}
\label{sec:family}

The gain is not uniform across families, and CGS is designed to exploit exactly this
structure. Figure~\ref{fig:family} and Table~\ref{tab:family} decompose the
rotations-to-target by family. On the dense \emph{molecular-like} and
\emph{transverse-field Ising} families the folded schedules cut the cost to reach
$0.99$ by roughly a third to a half and are the only schedules to reach $0.999$. On
\emph{Heisenberg}, by contrast, the commuting-clique structure makes the first-order
clique schedule essentially \emph{exact} at a single step
(Proposition~\ref{prop:colour}): every clique-based schedule reaches $0.999$ at the
minimal rotation count, and the symmetric fold's $2\times$ per-step overhead is here
wasted. The gate-optimal schedule is therefore genuinely instance dependent---fold
the hard instances, and leave the fast-forwardable ones first-order.

\begin{figure*}[t]
    \centering
    \includegraphics[width=0.9\textwidth]{fig_family.pdf}
    \caption{\textbf{Rotations to fidelity $\geq0.99$ by Hamiltonian family.} Lower is
    better; an $\times$ marks a target not reached within the swept budget. On the dense
    molecular-like and transverse-field Ising families the folded schedules
    (antithetic, symmetric) reach the target at substantially fewer rotations than any
    first-order schedule. On Heisenberg the commuting-clique first-order schedules are
    already near-exact at the minimal budget, so folding only adds overhead. Bars are
    over the $N=48$ instances per family.}
    \label{fig:family}
\end{figure*}

\begin{table}[t]
\centering
\caption{\textbf{Per-family rotations to a target fidelity.} For each family and
schedule, the realized rotation count to reach mean fidelity $\geq0.99$ and
$\geq0.999$; ``---'' marks a target not reached within the swept step grid. The folded
schedules are the only ones to reach $0.999$ on the molecular-like and transverse-field
Ising families; the clique-based first-order schedules already reach it at minimal cost
on Heisenberg. Transcribed from the \texttt{gates\_to\_target} block of
\texttt{results/summary.json}.}
\label{tab:family}
\input{tables/tab_family.tex}
\end{table}

The learned scheduler of Sec.~\ref{sec:router_thy} exploits this structure
end-to-end. Trained leave-one-instance-out over the ten commutator-graph features, so
that no test instance informs its own predictor, it attains $\LearnedAcc\%$ accuracy
against the per-instance oracle. We note that, on the subset of $\CommonN$ instances
where every fixed schedule reaches the target, the scheduler needs $\LearnedMeanGates$
rotations on average, against an oracle optimum of $\OracleMeanGates$ and the best
single fixed schedules' $\SymMeanGates$ (symmetric) and $\AntiMeanGates$
(antithetic)---recovering near-oracle efficiency by routing fast-forwardable instances
to the cheap first-order schedule and hard instances to a fold. The scheduler is
interpretable and listing-order invariant by construction
(Proposition~\ref{prop:invariance}): it is a selector over named schedules, not a
black-box policy.

% ============================================================
\section{Discussion}
\label{sec:discussion}

\paragraph{Summary of the result.}
In this manuscript we introduced commutator-graph scheduling (CGS), and the result it
establishes has two halves that must be kept together. \emph{Negatively}, a single
term ordering is a provably weak lever for first-order Trotterization: reordering only
re-signs the leading-error commutators, the leading-error norm is exactly invariant on
collision-free families and within a few percent of its floor otherwise, and no
ordering escapes the $O(t^2/r)$ rate---turning the well-known empirical observation
that ordering heuristics are within noise into a theorem. \emph{Positively}, the
freedom that does matter is the per-step schedule: the antithetic fold cancels the
leading error and converges at second order at the identical gate count, reaching a
target fidelity with $\GateSpeedup\times$ fewer rotations and reaching $0.999$ where
first-order Trotter plateaus. We do not claim a new product-formula
\emph{primitive}---the symmetrization mechanism is
classical~\cite{suzuki1991general,hatano2005finding}---but rather the framing that the
lever is the schedule, that a fixed ordering provably cannot realize it, and that the
commutator graph both certifies the impotence and structures the schedule, together
with the matched-cost quantification and the learned scheduler.

\paragraph{Relation to prior work.}
The commutator graph and its clique structure are established in measurement grouping
for variational
algorithms~\cite{verteletskyi2020measurement,gokhale2020optimization,crawford2021efficient};
we repurpose them for \emph{time evolution} and, beyond ordering, for scheduling. The
sensitivity of Trotterized chemistry to term ordering has been reported
empirically~\cite{tranter2019ordering,poulin2015trotter,hastings2015improving}, and the
impotence theorem explains the ceiling those studies encounter. Symmetric (Strang) and
higher-order product formulas are classical~\cite{strang1968construction,suzuki1990fractal,suzuki1991general,hatano2005finding},
and randomized compilation---random permutation per step, qDRIFT, stochastic
sparsification, randomized multi-product formulas, and composite channels---improves
error scaling by averaging rather than
cancellation~\cite{childs2019faster,campbell2019random,ouyang2020compilation,faehrmann2022randomizing,hagan2023composite};
the antithetic schedule is the
\emph{deterministic, zero-variance} counterpart that cancels the same leading term in a
single sample, and our matched-cost frontier places it against both the first-order and
the randomized baselines. A complementary deterministic route reduces the leading
error by taking linear combinations of Trotter circuits---multi-product
formulas~\cite{zhuk2024trotter}---at the cost of ancillas or postselection, whereas
the antithetic schedule pays no such overhead. Beyond product formulas, query-optimal
primitives based on linear combinations of unitaries, truncated Taylor series,
quantum signal processing, and qubitization achieve logarithmic precision
scaling~\cite{berry2007efficient,berry2015taylor,berry2015nearly,low2017optimal,low2019qubitization},
and refined input- and observable-dependent analyses tighten the cost
further~\cite{an2021time,kivlichan2017bounding}; these target the
fault-tolerant regime, while CGS is a near-term, zero-overhead upgrade. Hardware-aware
compilations such as the fermionic-swap network reduce circuit depth for the same
Trotter step~\cite{kivlichan2018quantum}. The tight commutator-scaling theory of Trotter
error~\cite{childs2021theory,childs2018toward} is the lens throughout, and the
commutator error form is its explicit, matrix-free, ordering-resolved representation.

\paragraph{Why a matched-cost schedule matters.}
For first-order formulas the schedule is genuinely free---it changes no gate, qubit, or
depth count---so any accuracy it recovers is pure profit. In the high-accuracy regime
where precise chemistry and materials simulation operate, the difference between an
$O(t^2/r)$ and an $O(t^3/r^2)$ method at fixed cost is the difference between reaching
and not reaching a target, as the unreachable $0.999$ column of Table~\ref{tab:gates}
shows. The commutator graph supplies both halves of the recommendation at no extra
cost: it certifies when ordering cannot help, and its clique coloring tells the learned
scheduler which instances to fold and which to leave alone.

\paragraph{Limitations.}
The boundaries of the study are explicit. (i)~The second-order advantage is a
small-step-size effect: at coarse steps the antithetic fold can be \emph{worse} than a
good first-order ordering (its curve in Fig.~\ref{fig:frontier} starts above and crosses
below near the intermediate budget); we report the crossover.
(ii)~The matrix-free Krylov reference extends the exact benchmark to $n=\MaxN$
qubits---and we verify there that the advantage is size independent
(Sec.~\ref{sec:scaling})---but truly large systems ($n\gtrsim30$), where even a
sparse statevector is intractable and a tensor-network reference is required, are not
reached; the exact \emph{spectral-norm} impotence sweep is itself capped at
$n\le\ImpMaxN$ (the matrix-free collision and Hilbert--Schmidt statistics carry
Theorem~\ref{thm:impotence} at all sizes).
(iii)~The families are transverse-field Ising, Heisenberg, and random local two-body
molecular-like sums, a proxy for, not an instance of, mapped molecular
electronic-structure Hamiltonians at
scale~\cite{bravyi2002fermionic,whitfield2011simulation}. (iv)~The cost is the realized
rotation count and ignores the gate-synthesis cost of individual rotations and hardware
connectivity; the antithetic schedule's seam rotations could be merged to reduce the
count further, which we do not exploit so that its budget matches first-order exactly.
(v)~All results are classical statevector simulations on CPU, with no hardware
validation.

\paragraph{Outlook.}
Most critically, because the antithetic upgrade changes no gate, qubit, or depth count,
commutator-graph scheduling constitutes a zero-overhead accuracy improvement that any
first-order Trotter compilation can adopt unchanged on near-term hardware. Having shown
that the advantage is size independent up to $n=\MaxN$ qubits with a matrix-free
reference, natural next steps are truly large systems with a tensor-network reference,
real mapped molecular Hamiltonians, higher-order folded schedules, and a graph neural
policy trained directly on a differentiable matched-cost objective. Two of these
directions are statistical at their core and worth stating as such. The learned
scheduler (Sec.~\ref{sec:router_thy}) is already a \emph{probabilistic surrogate}
for an expensive simulator: rather than run every schedule on every instance, it
predicts the gate-optimal choice from cheap, interpretable commutator-graph
features whose invariance we prove (Proposition~\ref{prop:invariance}), and the
same construction invites calibrated uncertainty on that prediction and
active-learning selection of which instances to simulate next---an experimental
design over training data that our leave-one-out protocol only begins to
exploit. In parallel, the shot-based certification of the residual error
(Theorem~\ref{thm:estimator}) is an estimation problem on hardware, and pushing it
to noisy devices---where the witness must be disentangled from device error by the
same estimation-plus-hypothesis-testing logic that underlies randomized
benchmarking---would let CGS's zero-overhead accuracy claim be validated
experimentally rather than by simulation alone. In short, by recasting the leading Trotter error
as a certifiable operator on the commutator graph, CGS reframes term ordering as the
weak lever it provably is, and identifies the per-step schedule as the free, structured
degree of freedom that simulation should exploit.

\appendix

\section{Proofs for the error-operator theory}
\label{app:proofs}
We collect the two proofs deferred from Sec.~\ref{sec:theory}: the quantitative
nested-commutator remainder (Lemma~\ref{lem:remainder}) and the full
$\mathsf{NP}$-hardness reduction (Theorem~\ref{thm:nphard}).

\subsection{Proof of Lemma~\ref{lem:remainder} (quantitative BCH/Magnus remainder)}
Write $A(\Delta)=\prod_{a=1}^{m}e^{-\ii\Delta B_a}$ with Hermitian $B_a$,
$\norm{B_a}=\abs{b_a}$ (a Pauli string has unit operator norm), and
$\mu=\sum_a\abs{b_a}$. Introduce the interpolation
$A(\Delta,s)=\prod_{a=1}^{m}e^{-\ii s\Delta B_a}$, $s\in[0,1]$, so $A(\Delta)=
A(\Delta,1)$, and let $A(\Delta,s)=\exp(\Omega(\Delta,s))$ be the (convergent, for
small $\Delta$) Magnus exponent with $\Omega(\Delta,0)=0$. Differentiating the
ordered product in $s$ and using the Magnus differential equation
$\dot\Omega=\sum_{k\ge0}\tfrac{B_k}{k!}\,\mathrm{ad}_\Omega^{\,k}\bigl(\dot A A^{-1}
\bigr)$ (with $B_k$ the Bernoulli numbers and $\mathrm{ad}_X(Y)=[X,Y]$), the
generator is the norm-convergent series
\[
    \Omega(\Delta,1)=\sum_{k\ge1}\Omega_k\Delta^k,
\]
where each $\Omega_k$ is a fixed linear combination of $k$-fold nested commutators
$[\,B_{a_1},[\,B_{a_2},[\dots,B_{a_k}]\dots]\,]$. Two elementary bounds control the
tail. First, each nested commutator of order $k$ satisfies, by the triangle
inequality $\norm{[X,Y]}\le2\norm{X}\,\norm{Y}$ applied $k-1$ times and
sub-multiplicativity across the ordered product,
\begin{equation*}
\begin{split}
    \Bigl\|\,k\text{-fold commutator of the }B_a\,\Bigr\|
    &\;\le\;(2)^{\,k-1}\Bigl(\sum_a\abs{b_a}\Bigr)^{k}\\
    &\;=\;2^{k-1}\mu^{k}.
\end{split}
\end{equation*}
Second, the absolute sum of the scalar Magnus/BCH coefficients at order $k$ is
dominated by the coefficients of the \emph{scalar} majorant series obtained by
replacing every $\norm{[X,Y]}$ by $2\norm{X}\norm{Y}$; that majorant is exactly the
series for $-\log(2-e^{2\mu\abs{\Delta}})$ about $\Delta=0$, which converges for
$2\mu\abs{\Delta}<\log2$. Truncating at order $p$ and setting
$G_p(\Delta)=\ii\sum_{k\le p}\Omega_k\Delta^k$ (Hermitian since each $\Omega_k$ is
anti-Hermitian, being a real combination of commutators of Hermitian operators times
$-\ii$), the remainder
$\Delta^{p+1}\mathcal{R}_p(\Delta)=\sum_{k>p}\Omega_k\Delta^k$ is bounded by the tail
of the majorant:
\begin{equation*}
\begin{split}
    \norm{\mathcal{R}_p(\Delta)}
    &\le\abs{\Delta}^{-(p+1)}\!\!\sum_{k>p}
    \frac{(2\mu\abs{\Delta})^{k}}{k}\\
    &\le\frac{(2\mu)^{p+1}}{p+1}\sum_{j\ge0}(2\mu\abs{\Delta})^{j}\\
    &=\frac{(2\mu)^{p+1}}{p+1}\cdot\frac{1}{1-2\mu\abs{\Delta}},
\end{split}
\end{equation*}
and $1/(1-2\mu\abs{\Delta})\le e^{2\mu\abs{\Delta}}$ for $2\mu\abs{\Delta}<\log2<1$,
giving Eq.~\eqref{eq:remainder}. The order-$1$ term is
$\Omega_1\Delta=-\ii\Delta\sum_a B_a$ (each factor contributes its generator
linearly), so $G_p$'s leading part is $\Delta\sum_aB_a$; the classical order-$2$ BCH
term of an ordered product is $-\tfrac{\Delta^2}{2}\sum_{a<b}[B_a,B_b]$ in the
exponent, i.e.\ the stated $-\tfrac{\Delta^2}{2}\sum_{a<b}\ii[B_a,B_b]$ inside
$-\ii G_p$. \hfill$\qed$

\subsection{Full reduction for Theorem~\ref{thm:nphard}}
Let $G=(V,E_G)$ be an instance of \textsc{Max-Cut}. We construct, in time polynomial
in $\abs{V}+\abs{E_G}$, a term list $\{(c_j,P_j)\}$ whose \textsc{MinResidual}
objective $F$ equals $2\abs{E_G}-2\sum_{uv\in E_G}x_ux_v$ under a bijection between
sign assignments and cuts, so that minimizing $F$ solves the (\textsc{NP}-hard)
minimum-agreement problem, equivalently maximizing the cut.

\emph{Gadget.} Assign to each vertex $u\in V$ a Boolean sign $x_u\in\{\pm1\}$, encoded
by the relative order of a designated ``anchor'' term $A$ and a per-vertex term $T_u$
in the schedule: $x_u=\varepsilon_{A,u}\in\{\pm1\}$ is exactly the sign the ordering
places on the anchor--$T_u$ edge (any acyclic tournament realizes an arbitrary such
sign pattern, so all $x\in\{\pm1\}^{V}$ are feasible). For each graph edge
$uv\in E_G$ introduce two CGS-edges that land on a \emph{shared} Pauli string
$R_{uv}$ (achievable by choosing $T_u,T_v,A$ so the two anticommuting products
$P_AP_{T_u}$ and $P_AP_{T_v}$ coincide up to phase---a collision, in the sense of
Definition~\ref{def:errorform}), with unit weights $w=1$ and base signs $s$ chosen so
that the group amplitude is $\alpha_{R_{uv}}=x_u-x_v$. Distinct graph edges use
disjoint collision groups, so
\[
    F(\varepsilon)=\sum_{uv\in E_G}\alpha_{R_{uv}}^2
    =\sum_{uv\in E_G}(x_u-x_v)^2 .
\]
\emph{Correctness.} Using $x_u^2=1$, $(x_u-x_v)^2=2-2x_ux_v\in\{0,4\}$: it is $0$ when
$x_u=x_v$ (endpoints uncut) and $4$ when $x_u\neq x_v$ (endpoints cut). Hence
$F=4\cdot\abs{\{uv:\text{cut}\}}$ and
$\tfrac14 F=\abs{E_G}-\tfrac12\sum_{uv}x_ux_v$. Therefore a sign assignment minimizing
$F$ minimizes the number of cut edges, equivalently \emph{maximizes} the agreement
$\sum_{uv}x_ux_v$; the complementary maximization of $F$ is \textsc{Max-Cut}. Because
computing \textsc{Max-Cut} (and its \textsc{Min-Uncut} complement) is
\textsc{NP}-hard and the reduction is polynomial, exactly minimizing $F$---i.e.\
solving \textsc{MinResidual}---is \textsc{NP}-hard. The construction never forms a
$2^n\times2^n$ matrix: each collision is verified by the $O(n)$ Pauli-product kernel
of Sec.~\ref{sec:pauli}, so the reduction is explicit and efficient. \hfill$\qed$

\section{Numerical methods and protocol}
\label{app:methods}
The analytic development---the Pauli algebra, the commutator error form, and the
ordering-impotence, antithetic-cancellation, clique-exactness, and
listing-order-invariance results---appears with complete proofs in
Sec.~\ref{sec:method}. Here we specify the numerical protocol. Three families are
generated as $\{(c_j,P_j)\}$ term lists: transverse-field Ising
$H=-J\sum_i Z_iZ_{i+1}-h\sum_i X_i$; Heisenberg
$H=\sum_i(J_x X_iX_{i+1}+J_y Y_iY_{i+1}+J_z Z_iZ_{i+1})$; and \emph{molecular-like},
random local two-body Pauli terms with seeded Gaussian coefficients---a tractable
stand-in for the dense, irregular anticommutation structure of Jordan--Wigner /
Bravyi--Kitaev mapped electronic-structure
Hamiltonians~\cite{bravyi2002fermionic,whitfield2011simulation,jordan1928paulische,seeley2012bravyi,tranter2015bravyi,babbush2016exponentially}. The reported-scale
configuration (\texttt{configs/full.yaml}) uses sizes $n\in\{4,6,8,10,12,14\}$, eight
seeded instances per (family, size), evolution time $t=\EvolTime$, a reference budget
of $r=\RefSteps$ steps, and a step grid $\{1,2,4,6,8,12,16,24\}$ for the convergence
and frontier curves, giving $\Ninstances$ instances in total. The exact reference
$e^{-\ii Ht}\ket{\psi_0}$ is computed either by dense \texttt{scipy.linalg.expm} (the
ground truth, used for $n\le\ImpMaxN$) or by the \emph{matrix-free} Krylov reference
\texttt{expm\_multiply} applied to a sparse Hamiltonian assembled term-by-term from
one-nonzero-per-column Pauli permutations [Eq.~\eqref{eq:rotation} extended to the
reference]; the latter never forms a $2^n\times2^n$ matrix, agrees with dense
\texttt{expm} to $10^{-9}$ (the \texttt{reference\_backend\_verified} flag), and is
what extends the exact benchmark to $n=\MaxN$. For each instance and schedule we
evaluate the product formula, score fidelity and infidelity against the reference, and
record the realized rotation count; from the (rotation count, fidelity) curves we read
the rotations to each target fidelity and fit the convergence rate. Ordering impotence
is measured by sampling $48$ orderings per instance and forming the exact leading-error
spectral norm on the $\ImpNInst$ instances with $n\le\ImpMaxN$, where the dense
$O((2^n)^3)$ spectral norm is tractable; the matrix-free Hilbert--Schmidt surrogate and
collision statistics are computed at all sizes. Per-schedule means carry 95\%
confidence intervals $1.96\,\hat\sigma/\sqrt{N}$.

\section{Software, validation, and reproducibility}
\label{app:repro}
The pipeline is a CPU-only Python package, distributed on pip as
\texttt{specops-cgs} (import name \texttt{topoham}), depending on
NumPy~\cite{harris2020array}, SciPy~\cite{virtanen2020scipy},
NetworkX~\cite{hagberg2008exploring}, scikit-learn~\cite{pedregosa2011scikit}, and
Matplotlib. Hamiltonians are stored as Pauli term lists (\texttt{hamiltonians.py}); the
Pauli algebra of Sec.~\ref{sec:pauli}---the $O(n)$ parity test (Lemma~\ref{lem:parity}),
the cosine--sine rotation kernel (Lemma~\ref{lem:rotation}), and the phased-Pauli
product---is the inner loop of the exact statevector simulator (\texttt{pauli.py}). The
commutator graph and matrix-free error form are built in \texttt{commutator\_graph.py}
and \texttt{error\_form.py}: the greedy coloring partitions terms into commuting cliques
(Proposition~\ref{prop:colour}), the error form assembles Eq.~\eqref{eq:edgesum} as a
signed Pauli-string sum over the edges (Definition~\ref{def:errorform}), and the same
module computes the ten listing-order-invariant features the learned scheduler reads
(Proposition~\ref{prop:invariance}). Schedules are realized as flat lists of
(term, step-fraction) rotations whose seam-merged length is the gate cost on which all
schedules are compared (\texttt{schedules.py}, \texttt{env.py}); the learned scheduler
(\texttt{LearnedSchedulePolicy}) is the standardized logistic-regression classifier of
Sec.~\ref{sec:router_thy}, trained leave-one-instance-out. The exact reference is built
in \texttt{env.py}: \texttt{sparse\_pauli\_matrix} assembles each Pauli string as a
one-nonzero-per-column signed permutation directly from index arithmetic in $O(2^n)$,
so the sparse Hamiltonian and its \texttt{expm\_multiply} reference are formed without
a dense $2^n\times2^n$ matrix.

The single source of truth is \texttt{results/summary.json}, written by the runner with
full provenance (seed, command, platform, package versions, runtime, peak memory) and
\emph{three} integrity flags that must pass before any fidelity is produced:
\texttt{pauli\_algebra\_verified} (the fast kernel matches dense \texttt{expm} to
$10^{-9}$), \texttt{error\_form\_verified} (the matrix-free error form matches the dense
pair-commutator sum to $10^{-9}$), and \texttt{reference\_backend\_verified} (the
matrix-free Krylov reference matches dense \texttt{expm} to $10^{-9}$ before it is used
at scale). On the run reported here (seed $0$, Python~3.9, macOS arm64), the experiment
completed in $\approx\Runtime$\,s wall-clock with $\approx\PeakMem$\,MB peak memory, all
three flags \texttt{true}. The macros, tables,
and figures in this manuscript are generated directly from \texttt{summary.json}, so no
number is entered by hand. The full pipeline is reproduced from the accompanying code
(\texttt{code/} in the submission; package \texttt{specops-cgs}) by
\begin{verbatim}
pip install .
make test
export KMP_DUPLICATE_LIB_OK=TRUE
cgs-reproduce \
  --config configs/full.yaml
\end{verbatim}
where \texttt{make test} checks the Pauli kernel ($10^{-9}$), the error form
($10^{-9}$), antithetic second-order convergence, and clique exactness, and
\texttt{cgs-reproduce} regenerates \texttt{summary.json}, the tables, and the
figures, then syncs them into \texttt{figures/} and \texttt{tables/} used by this
manuscript; \texttt{cgs-reproduce --quick} (equivalently
\texttt{--config configs/smoke.yaml}) runs a seconds-scale demonstration of the
same pipeline. All experiments are reproducible on commodity hardware.

\begin{thebibliography}{99}
\providecommand{\natexlab}[1]{#1}
\providecommand{\url}[1]{\texttt{#1}}
\expandafter\ifx\csname urlstyle\endcsname\relax
  \providecommand{\doi}[1]{doi: #1}\else
  \providecommand{\doi}{doi: \begingroup \urlstyle{rm}\Url}\fi

\bibitem{feynman1982simulating}
Richard~P. Feynman.
\newblock Simulating physics with computers.
\newblock \emph{International Journal of Theoretical Physics}, 21:\penalty0
  467--488, 1982.
\newblock \doi{10.1007/BF02650179}.

\bibitem{lloyd1996universal}
Seth Lloyd.
\newblock Universal quantum simulators.
\newblock \emph{Science}, 273\penalty0 (5278):\penalty0 1073--1078, 1996.
\newblock \doi{10.1126/science.273.5278.1073}.

\bibitem{trotter1959product}
Hale~F. Trotter.
\newblock On the product of semi-groups of operators.
\newblock \emph{Proceedings of the American Mathematical Society}, 10\penalty0
  (4):\penalty0 545--551, 1959.
\newblock \doi{10.1090/S0002-9939-1959-0108732-6}.

\bibitem{suzuki1991general}
Masuo Suzuki.
\newblock General theory of fractal path integrals with applications to
  many-body theories and statistical physics.
\newblock \emph{Journal of Mathematical Physics}, 32\penalty0 (2):\penalty0
  400--407, 1991.
\newblock \doi{10.1063/1.529425}.

\bibitem{childs2021theory}
Andrew~M. Childs, Yuan Su, Minh~C. Tran, Nathan Wiebe, and Shuchen Zhu.
\newblock Theory of trotter error with commutator scaling.
\newblock \emph{Physical Review X}, 11\penalty0 (1):\penalty0 011020, 2021.
\newblock \doi{10.1103/PhysRevX.11.011020}.

\bibitem{childs2018toward}
Andrew~M. Childs, Dmitri Maslov, Yunseong Nam, Neil~J. Ross, and Yuan Su.
\newblock Toward the first quantum simulation with quantum speedup.
\newblock \emph{Proceedings of the National Academy of Sciences}, 115\penalty0
  (38):\penalty0 9456--9461, 2018.
\newblock \doi{10.1073/pnas.1801723115}.

\bibitem{hatano2005finding}
Naomichi Hatano and Masuo Suzuki.
\newblock Finding exponential product formulas of higher orders.
\newblock \emph{Lecture Notes in Physics}, 679:\penalty0 37--68, 2005.
\newblock \doi{10.1007/11526216_2}.

\bibitem{childs2019faster}
Andrew~M. Childs, Aaron Ostrander, and Yuan Su.
\newblock Faster quantum simulation by randomization.
\newblock \emph{Quantum}, 3:\penalty0 182, 2019.
\newblock \doi{10.22331/q-2019-09-02-182}.

\bibitem{campbell2019random}
Earl Campbell.
\newblock Random compiler for fast Hamiltonian simulation.
\newblock \emph{Physical Review Letters}, 123\penalty0 (7):\penalty0 070503,
  2019.
\newblock \doi{10.1103/PhysRevLett.123.070503}.

\bibitem{poulin2015trotter}
David Poulin, Matthew~B. Hastings, Dave Wecker, Nathan Wiebe, Andrew~C.
  Doherty, and Matthias Troyer.
\newblock The trotter step size required for accurate quantum simulation of
  quantum chemistry.
\newblock \emph{Quantum Information \& Computation}, 15\penalty0
  (5--6):\penalty0 361--384, 2015.

\bibitem{hastings2015improving}
Matthew~B. Hastings, Dave Wecker, Bela Bauer, and Matthias Troyer.
\newblock Improving quantum algorithms for quantum chemistry.
\newblock \emph{Quantum Information \& Computation}, 15\penalty0
  (1--2):\penalty0 1--21, 2015.

\bibitem{verteletskyi2020measurement}
Vladyslav Verteletskyi, Tzu-Ching Yen, and Artur~F. Izmaylov.
\newblock Measurement optimization in the variational quantum eigensolver using
  a minimum clique cover.
\newblock \emph{The Journal of Chemical Physics}, 152\penalty0 (12):\penalty0
  124114, 2020.
\newblock \doi{10.1063/1.5141458}.

\bibitem{gokhale2020optimization}
Pranav Gokhale, Olivia Angiuli, Yongshan Ding, Kaiwen Gui, Teague Tomesh,
  Martin Suchara, Margaret Martonosi, and Frederic~T. Chong.
\newblock {$O(N^3)$} measurement cost for variational quantum eigensolver on
  molecular hamiltonians.
\newblock \emph{IEEE Transactions on Quantum Engineering}, 1:\penalty0 1--24,
  2020.
\newblock \doi{10.1109/TQE.2020.3035814}.

\bibitem{crawford2021efficient}
Ophelia Crawford, Billy~van Straaten, Daochen Wang, Thomas Parks, Earl
  Campbell, and Stephen Brierley.
\newblock Efficient quantum measurement of pauli operators in the presence of
  finite sampling error.
\newblock \emph{Quantum}, 5:\penalty0 385, 2021.
\newblock \doi{10.22331/q-2021-01-20-385}.

\bibitem{tranter2019ordering}
Andrew Tranter, Peter~J. Love, Florian Mintert, and Nathan Wiebe.
\newblock Ordering of trotterization: Impact on errors in quantum simulation of
  electronic structure.
\newblock \emph{Entropy}, 21\penalty0 (12):\penalty0 1218, 2019.
\newblock \doi{10.3390/e21121218}.

\bibitem{pedregosa2011scikit}
Fabian Pedregosa, Ga{\"e}l Varoquaux, Alexandre Gramfort, et~al.
\newblock Scikit-learn: Machine learning in python.
\newblock \emph{Journal of Machine Learning Research}, 12:\penalty0 2825--2830,
  2011.

\bibitem{bravyi2002fermionic}
Sergey~B. Bravyi and Alexei~Yu. Kitaev.
\newblock Fermionic quantum computation.
\newblock \emph{Annals of Physics}, 298\penalty0 (1):\penalty0 210--226, 2002.
\newblock \doi{10.1006/aphy.2002.6254}.

\bibitem{whitfield2011simulation}
James~D. Whitfield, Jacob Biamonte, and Al{\'a}n Aspuru-Guzik.
\newblock Simulation of electronic structure hamiltonians using quantum
  computers.
\newblock \emph{Molecular Physics}, 109\penalty0 (5):\penalty0 735--750, 2011.
\newblock \doi{10.1080/00268976.2011.552441}.

\bibitem{harris2020array}
Charles~R. Harris, K.~Jarrod Millman, St{\'e}fan~J. van~der Walt, et~al.
\newblock Array programming with numpy.
\newblock \emph{Nature}, 585:\penalty0 357--362, 2020.
\newblock \doi{10.1038/s41586-020-2649-2}.

\bibitem{virtanen2020scipy}
Pauli Virtanen, Ralf Gommers, Travis~E. Oliphant, et~al.
\newblock Scipy 1.0: fundamental algorithms for scientific computing in python.
\newblock \emph{Nature Methods}, 17:\penalty0 261--272, 2020.
\newblock \doi{10.1038/s41592-019-0686-2}.

\bibitem{hagberg2008exploring}
Aric~A. Hagberg, Daniel~A. Schult, and Pieter~J. Swart.
\newblock Exploring network structure, dynamics, and function using networkx.
\newblock In \emph{Proceedings of the 7th Python in Science Conference}, pages
  11--15, 2008.

\bibitem{georgescu2014quantum}
I.~M. Georgescu, S.~Ashhab, and Franco Nori.
\newblock Quantum simulation.
\newblock \emph{Reviews of Modern Physics}, 86\penalty0 (1):\penalty0 153--185,
  2014.
\newblock \doi{10.1103/RevModPhys.86.153}.

\bibitem{aspuru2005simulated}
Al{\'a}n Aspuru-Guzik, Anthony~D. Dutoi, Peter~J. Love, and Martin Head-Gordon.
\newblock Simulated quantum computation of molecular energies.
\newblock \emph{Science}, 309\penalty0 (5741):\penalty0 1704--1707, 2005.
\newblock \doi{10.1126/science.1113479}.

\bibitem{mcardle2020quantum}
Sam McArdle, Suguru Endo, Al{\'a}n Aspuru-Guzik, Simon~C. Benjamin, and Xiao
  Yuan.
\newblock Quantum computational chemistry.
\newblock \emph{Reviews of Modern Physics}, 92\penalty0 (1):\penalty0 015003,
  2020.
\newblock \doi{10.1103/RevModPhys.92.015003}.

\bibitem{lanyon2011universal}
B.~P. Lanyon, C.~Hempel, D.~Nigg, M.~M{\"u}ller, R.~Gerritsma, F.~Z{\"a}hringer,
  P.~Schindler, J.~T. Barreiro, M.~Rambach, G.~Kirchmair, M.~Hennrich,
  P.~Zoller, R.~Blatt, and C.~F. Roos.
\newblock Universal digital quantum simulation with trapped ions.
\newblock \emph{Science}, 334\penalty0 (6052):\penalty0 57--61, 2011.
\newblock \doi{10.1126/science.1208001}.

\bibitem{smith2016manybody}
J.~Smith, A.~Lee, P.~Richerme, B.~Neyenhuis, P.~W. Hess, P.~Hauke, M.~Heyl,
  D.~A. Huse, and C.~Monroe.
\newblock Many-body localization in a quantum simulator with programmable random
  disorder.
\newblock \emph{Nature Physics}, 12\penalty0 (10):\penalty0 907--911, 2016.
\newblock \doi{10.1038/nphys3783}.

\bibitem{strang1968construction}
Gilbert Strang.
\newblock On the construction and comparison of difference schemes.
\newblock \emph{SIAM Journal on Numerical Analysis}, 5\penalty0 (3):\penalty0
  506--517, 1968.
\newblock \doi{10.1137/0705041}.

\bibitem{suzuki1990fractal}
Masuo Suzuki.
\newblock Fractal decomposition of exponential operators with applications to
  many-body theories and Monte Carlo simulations.
\newblock \emph{Physics Letters A}, 146\penalty0 (6):\penalty0 319--323, 1990.
\newblock \doi{10.1016/0375-9601(90)90962-N}.

\bibitem{berry2007efficient}
Dominic~W. Berry, Graeme Ahokas, Richard Cleve, and Barry~C. Sanders.
\newblock Efficient quantum algorithms for simulating sparse Hamiltonians.
\newblock \emph{Communications in Mathematical Physics}, 270\penalty0
  (2):\penalty0 359--371, 2007.
\newblock \doi{10.1007/s00220-006-0150-x}.

\bibitem{childs2012hamiltonian}
Andrew~M. Childs and Nathan Wiebe.
\newblock Hamiltonian simulation using linear combinations of unitary
  operations.
\newblock \emph{Quantum Information \& Computation}, 12\penalty0
  (11--12):\penalty0 901--924, 2012.

\bibitem{childs2010limitations}
Andrew~M. Childs and Robin Kothari.
\newblock Limitations on the simulation of non-sparse Hamiltonians.
\newblock \emph{Quantum Information \& Computation}, 10\penalty0
  (7--8):\penalty0 669--684, 2010.

\bibitem{berry2015taylor}
Dominic~W. Berry, Andrew~M. Childs, Richard Cleve, Robin Kothari, and Rolando~D.
  Somma.
\newblock Simulating Hamiltonian dynamics with a truncated Taylor series.
\newblock \emph{Physical Review Letters}, 114\penalty0 (9):\penalty0 090502,
  2015.
\newblock \doi{10.1103/PhysRevLett.114.090502}.

\bibitem{berry2015nearly}
Dominic~W. Berry, Andrew~M. Childs, and Robin Kothari.
\newblock Hamiltonian simulation with nearly optimal dependence on all
  parameters.
\newblock In \emph{2015 IEEE 56th Annual Symposium on Foundations of Computer
  Science (FOCS)}, pages 792--809. IEEE, 2015.
\newblock \doi{10.1109/FOCS.2015.54}.

\bibitem{low2017optimal}
Guang~Hao Low and Isaac~L. Chuang.
\newblock Optimal Hamiltonian simulation by quantum signal processing.
\newblock \emph{Physical Review Letters}, 118\penalty0 (1):\penalty0 010501,
  2017.
\newblock \doi{10.1103/PhysRevLett.118.010501}.

\bibitem{low2019qubitization}
Guang~Hao Low and Isaac~L. Chuang.
\newblock Hamiltonian simulation by qubitization.
\newblock \emph{Quantum}, 3:\penalty0 163, 2019.
\newblock \doi{10.22331/q-2019-07-12-163}.

\bibitem{childs2019lattice}
Andrew~M. Childs and Yuan Su.
\newblock Nearly optimal lattice simulation by product formulas.
\newblock \emph{Physical Review Letters}, 123\penalty0 (5):\penalty0 050503,
  2019.
\newblock \doi{10.1103/PhysRevLett.123.050503}.

\bibitem{tran2020destructive}
Minh~C. Tran, Su-Kuan Chu, Yuan Su, Andrew~M. Childs, and Alexey~V. Gorshkov.
\newblock Destructive error interference in product-formula lattice simulation.
\newblock \emph{Physical Review Letters}, 124\penalty0 (22):\penalty0 220502,
  2020.
\newblock \doi{10.1103/PhysRevLett.124.220502}.

\bibitem{layden2022first}
David Layden.
\newblock First-order Trotter error from a second-order perspective.
\newblock \emph{Physical Review Letters}, 128\penalty0 (21):\penalty0 210501,
  2022.
\newblock \doi{10.1103/PhysRevLett.128.210501}.

\bibitem{yi2022spectral}
Changhao Yi and Elizabeth Crosson.
\newblock Spectral analysis of product formulas for quantum simulation.
\newblock \emph{npj Quantum Information}, 8\penalty0 (1):\penalty0 37, 2022.
\newblock \doi{10.1038/s41534-022-00548-w}.

\bibitem{chen2024averagecase}
Chi-Fang Chen and Fernando G.~S.~L. Brand{\~a}o.
\newblock Average-case speedup for product formulas.
\newblock \emph{Communications in Mathematical Physics}, 405\penalty0
  (2):\penalty0 32, 2024.
\newblock \doi{10.1007/s00220-023-04912-5}.

\bibitem{heyl2019quantum}
Markus Heyl, Philipp Hauke, and Peter Zoller.
\newblock Quantum localization bounds Trotter errors in digital quantum
  simulation.
\newblock \emph{Science Advances}, 5\penalty0 (4):\penalty0 eaau8342, 2019.
\newblock \doi{10.1126/sciadv.aau8342}.

\bibitem{sieberer2019digital}
L.~M. Sieberer, T.~Olsacher, A.~Elben, M.~Heyl, P.~Hauke, F.~Haake, and
  P.~Zoller.
\newblock Digital quantum simulation, Trotter errors, and quantum chaos of the
  kicked top.
\newblock \emph{npj Quantum Information}, 5\penalty0 (1):\penalty0 78, 2019.
\newblock \doi{10.1038/s41534-019-0192-5}.

\bibitem{ouyang2020compilation}
Yingkai Ouyang, David~R. White, and Earl~T. Campbell.
\newblock Compilation by stochastic Hamiltonian sparsification.
\newblock \emph{Quantum}, 4:\penalty0 235, 2020.
\newblock \doi{10.22331/q-2020-02-27-235}.

\bibitem{faehrmann2022randomizing}
Paul~K. Faehrmann, Mark Steudtner, Richard Kueng, Maria Kiefer{\'o}v{\'a}, and
  Jens Eisert.
\newblock Randomizing multi-product formulas for Hamiltonian simulation.
\newblock \emph{Quantum}, 6:\penalty0 806, 2022.
\newblock \doi{10.22331/q-2022-09-19-806}.

\bibitem{zhuk2024trotter}
Sergiy Zhuk, Niall Robertson, and Sergey Bravyi.
\newblock Trotter error bounds and dynamic multi-product formulas for
  Hamiltonian simulation.
\newblock \emph{Physical Review Research}, 6\penalty0 (3):\penalty0 033309,
  2024.
\newblock \doi{10.1103/PhysRevResearch.6.033309}.

\bibitem{hagan2023composite}
Matthew Hagan and Nathan Wiebe.
\newblock Composite quantum simulations.
\newblock \emph{Quantum}, 7:\penalty0 1181, 2023.
\newblock \doi{10.22331/q-2023-11-14-1181}.

\bibitem{an2021time}
Dong An, Di~Fang, and Lin Lin.
\newblock Time-dependent unbounded Hamiltonian simulation with vector norm
  scaling.
\newblock \emph{Quantum}, 5:\penalty0 459, 2021.
\newblock \doi{10.22331/q-2021-05-26-459}.

\bibitem{wecker2015solving}
Dave Wecker, Matthew~B. Hastings, Nathan Wiebe, Bryan~K. Clark, Chetan Nayak,
  and Matthias Troyer.
\newblock Solving strongly correlated electron models on a quantum computer.
\newblock \emph{Physical Review A}, 92\penalty0 (6):\penalty0 062318, 2015.
\newblock \doi{10.1103/PhysRevA.92.062318}.

\bibitem{babbush2016exponentially}
Ryan Babbush, Dominic~W. Berry, Ian~D. Kivlichan, Annie~Y. Wei, Peter~J. Love,
  and Al{\'a}n Aspuru-Guzik.
\newblock Exponentially more precise quantum simulation of fermions in second
  quantization.
\newblock \emph{New Journal of Physics}, 18\penalty0 (3):\penalty0 033032,
  2016.
\newblock \doi{10.1088/1367-2630/18/3/033032}.

\bibitem{kivlichan2018quantum}
Ian~D. Kivlichan, Jarrod McClean, Nathan Wiebe, Craig Gidney, Al{\'a}n
  Aspuru-Guzik, Garnet Kin-Lic Chan, and Ryan Babbush.
\newblock Quantum simulation of electronic structure with linear depth and
  connectivity.
\newblock \emph{Physical Review Letters}, 120\penalty0 (11):\penalty0 110501,
  2018.
\newblock \doi{10.1103/PhysRevLett.120.110501}.

\bibitem{kivlichan2017bounding}
Ian~D. Kivlichan, Nathan Wiebe, Ryan Babbush, and Al{\'a}n Aspuru-Guzik.
\newblock Bounding the costs of quantum simulation of many-body physics in real
  space.
\newblock \emph{Journal of Physics A: Mathematical and Theoretical},
  50\penalty0 (30):\penalty0 305301, 2017.
\newblock \doi{10.1088/1751-8121/aa77b8}.

\bibitem{jordan1928paulische}
Pascual Jordan and Eugene Wigner.
\newblock {\"U}ber das Paulische {\"A}quivalenzverbot.
\newblock \emph{Zeitschrift f{\"u}r Physik}, 47\penalty0 (9--10):\penalty0
  631--651, 1928.
\newblock \doi{10.1007/BF01331938}.

\bibitem{seeley2012bravyi}
Jacob~T. Seeley, Martin~J. Richard, and Peter~J. Love.
\newblock The Bravyi-Kitaev transformation for quantum computation of electronic
  structure.
\newblock \emph{The Journal of Chemical Physics}, 137\penalty0 (22):\penalty0
  224109, 2012.
\newblock \doi{10.1063/1.4768229}.

\bibitem{tranter2015bravyi}
Andrew Tranter, Sarah Sofia, Jake Seeley, Michael Kaicher, Jarrod McClean, Ryan
  Babbush, Peter~V. Coveney, Florian Mintert, Frank Wilhelm, and Peter~J. Love.
\newblock The Bravyi-Kitaev transformation: Properties and applications.
\newblock \emph{International Journal of Quantum Chemistry}, 115\penalty0
  (19):\penalty0 1431--1441, 2015.
\newblock \doi{10.1002/qua.24969}.

\bibitem{yen2020measuring}
Tzu-Ching Yen, Vladyslav Verteletskyi, and Artur~F. Izmaylov.
\newblock Measuring all compatible operators in one series of single-qubit
  measurements using unitary transformations.
\newblock \emph{Journal of Chemical Theory and Computation}, 16\penalty0
  (4):\penalty0 2400--2409, 2020.
\newblock \doi{10.1021/acs.jctc.0c00008}.

\bibitem{jena2019pauli}
Andrew Jena, Scott Genin, and Michele Mosca.
\newblock Pauli partitioning with respect to gate sets.
\newblock \emph{arXiv preprint arXiv:1907.07859}, 2019.

\bibitem{huynh2026gctr} M.~Huynh. Query-Efficient Quantum Approximate Optimization via Graph-Conditioned Trust Regions. arXiv:2604.24803 (2026).
\bibitem{huynh2026gctrcalpro} M.~Huynh. Graph-conditioned trust regions for uncertainty-calibrated quantum approximate optimization (extended). Manuscript in preparation (2026).
\bibitem{huynh2026meascost} M.~Huynh. The measurement cost of warm-started low-depth QAOA. Manuscript in preparation (2026).
\bibitem{huynh2026certqb} M.~Huynh. Certified query budgets for the quantum approximate optimization algorithm. Manuscript in preparation (2026).
\bibitem{huynh2026paramxferB} M.~Huynh. Topology-Conditioned QAOA Parameter Transfer for Budgeted Graph Optimization (companion). Manuscript in preparation (2026).
\bibitem{huynh2026ostqaoa} M.~Huynh. Operator-Spectral Truncated Priors for Query-Efficient QAOA Parameter Search. Manuscript in preparation (2026).
\bibitem{huynh2026topoqaoa} M.~Huynh. Spectral-truncation graph kernels for QAOA warm starts: topology-conditioned schedule transfer beyond depth one. Manuscript in preparation (2026).

\end{thebibliography}

\end{document}
